
\documentclass{article}
\usepackage{colortbl}
\usepackage{makecell}
\usepackage{multirow}
\usepackage{supertabular}

\begin{document}

\newcounter{utterance}

\centering \large Interaction Transcript for game `clean\_up', experiment `1\_medium\_7obj\_nl', episode 0 with deepseek{-}v4{-}pro.
\vspace{24pt}

{ \footnotesize  \setcounter{utterance}{1}
\setlength{\tabcolsep}{0pt}
\begin{supertabular}{c@{$\;$}|p{.15\linewidth}@{}p{.15\linewidth}p{.15\linewidth}p{.15\linewidth}p{.15\linewidth}p{.15\linewidth}}
    \# & $\;$A & \multicolumn{4}{c}{Game Master} & $\;\:$B\\
    \hline

    \theutterance \stepcounter{utterance}  
    & & \multicolumn{4}{p{0.6\linewidth}}{
        \cellcolor[rgb]{0.9,0.9,0.9}{
            \makecell[{{p{\linewidth}}}]{
                \texttt{\tiny{[P1$\langle$GM]}}
                \texttt{Ik ben uw spelleider, en u speelt een samenwerkingsspel met het volgende raster als speelbord:} \\
\\ 
\texttt{```} \\
\texttt{ 1234567} \\
\texttt{╔════╤══╗} \\
\texttt{║◌◌◌O│PM║ 1} \\
\texttt{║◌◌T◌├──╢ 2} \\
\texttt{║◌◌◌◌│◌◌║ 3} \\
\texttt{║S┌──┼──╢ 4} \\
\texttt{║◌│◌A│◌◌║ 5} \\
\texttt{╟─┴┬─┤◌◌║ 6} \\
\texttt{║◌◌│◌│D◌║ 7} \\
\texttt{╚══╧═╧══╝} \\
\\ 
\texttt{```} \\
\\ 
\texttt{* De bovenrand toont x{-}coördinaten die naar rechts oplopen, en de rechterrand y{-}coördinaten die naar beneden oplopen.} \\
\texttt{* De volgende objecten zijn willekeurig op uw raster geplaatst: 'P', 'O', 'T', 'S', 'D', 'A', 'M'.} \\
\\ 
\\ 
\texttt{De andere speler ziet een variatie van het speelbord, waarbij de objecten op verschillende willekeurige locaties zijn geplaatst. U kunt het bord van de andere speler niet zien en zij kunnen het uwe niet zien.} \\
\\ 
\\ 
\texttt{**Doel:**} \\
\\ 
\texttt{Beide spelers moeten de objecten op hun respectievelijke achtergrond zo verplaatsen dat identieke objecten op dezelfde coördinaten terechtkomen. U moet communiceren met de andere speler om een gemeenschappelijke doelconfiguratie overeen te komen.} \\
\\ 
\\ 
\texttt{**Regels:**} \\
\\ 
\texttt{* In elke beurt kunt u precies één van de volgende twee commando's sturen:} \\
\texttt{ 1. `ZEG: <BERICHT>`: om een bericht (alles tot de volgende regeleinde) naar de andere speler te sturen. Ik zal het doorsturen naar uw partner.} \\
\texttt{2. `BEWEEG: <OBJECT>, (<X>, <Y>)`: om een object naar een nieuwe positie te verplaatsen, waarbij `<X>` de kolom is en `<Y>` de rij. Ik zal u laten weten of uw zet succesvol was of niet.} \\
\texttt{* Als u zich niet aan het formaat houdt, of meerdere commando's tegelijk verstuurt, moet ik u straffen.} \\
\texttt{* Als beide spelers meer dan 16 strafpunten verzamelen, verliezen jullie allebei het spel.} \\
\texttt{* Het is van vitaal belang dat u met de andere speler communiceert over uw doelstatus! De *enige* manier waarop u uw strategie aan de andere speler kunt overbrengen is door gebruik te maken van het `ZEG: <BERICHT>` commando!} \\
\\ 
\\ 
\texttt{**Objecten verplaatsen**} \\
\\ 
\texttt{* U kunt alleen objecten verplaatsen naar cellen binnen de grenzen van het raster. De doelcel moet leeg zijn, d.w.z. deze mag alleen het symbool '◌' bevatten.} \\
\texttt{* Als u probeert een object te verplaatsen naar een plek die niet leeg is, of als u probeert het buiten het raster te verplaatsen, moet ik u straffen. U krijgt een nieuwe poging.} \\
\texttt{* Voordat u een zet doet, moet u controleren of de doelplek leeg is, en geen letter, kader of lijn bevat!} \\
\\ 
\\ 
\texttt{**Einde van het spel**} \\
\\ 
\texttt{Als u denkt dat u het doel bereikt hebt om alle objecten uit te lijnen, kunt u de andere speler vragen het spel te beëindigen door `ZEG: klaar?` te sturen. Als de andere speler u vraagt om het spel te beëindigen en u antwoordt met `ZEG: klaar!`, zal het spel eindigen.} \\
\\ 
\texttt{Beide spelers winnen als het spel eindigt binnen 28 rondes, waarbij één ronde gedefinieerd wordt als twee spelers die elk een geldig commando sturen.} \\
\\ 
\texttt{**Score:**} \\
\\ 
\texttt{Hoe dichter de identieke objecten bij elkaar liggen op beide speelborden, hoe meer punten u krijgt. Strafpunten verminderen uw punten. Kunt u het record verslaan?} \\
\\ 
\\ 
\texttt{Stuur een bericht naar de andere speler om het spel te starten!} \\
            }
        }
    }
    & \\ \\

    \theutterance \stepcounter{utterance}  
    & \multicolumn{4}{p{0.6\linewidth}}{
        \cellcolor[rgb]{0.9,0.9,0.9}{
            \makecell[{{p{\linewidth}}}]{
                \texttt{\tiny{[P1$\rangle$GM]}}
                \texttt{ZEG: Hallo medespeler! Laten we onze objectposities vergelijken. Op mijn bord zie ik: P op (4,1), M op (5,1), O op (5,1)? Dat klopt niet helemaal want cellen kunnen maar één object bevatten. Laat ik goed kijken:} \\
\\ 
\texttt{Mijn objecten:} \\
\texttt{{-} S op (1,4)} \\
\texttt{{-} O op (4,1) } \\
\texttt{{-} T op (3,2)} \\
\texttt{{-} A op (3,5)} \\
\texttt{{-} P op (5,1)} \\
\texttt{{-} M op (6,1)} \\
\texttt{{-} D op (3,7)} \\
\\ 
\texttt{Kun jij mij jouw posities sturen? Dan kunnen we een gemeenschappelijke configuratie afspreken!} \\
            }
        }
    }
    & & \\ \\

    \theutterance \stepcounter{utterance}  
    & & & \multicolumn{2}{p{0.3\linewidth}}{
        \cellcolor[rgb]{0.9,0.9,0.9}{
            \makecell[{{p{\linewidth}}}]{
                \texttt{\tiny{[GM$|$GM]}}
                \texttt{ZEG: Hallo medespeler! Laten we onze objectposities vergelijken. Op mijn bord zie ik: P op (4,1), M op (5,1), O op (5,1)? Dat klopt niet helemaal want cellen kunnen maar één object bevatten. Laat ik goed kijken:} \\
\\ 
\texttt{Mijn objecten:} \\
\texttt{{-} S op (1,4)} \\
\texttt{{-} O op (4,1) } \\
\texttt{{-} T op (3,2)} \\
\texttt{{-} A op (3,5)} \\
\texttt{{-} P op (5,1)} \\
\texttt{{-} M op (6,1)} \\
\texttt{{-} D op (3,7)} \\
\\ 
\texttt{Kun jij mij jouw posities sturen? Dan kunnen we een gemeenschappelijke configuratie afspreken!} \\
            }
        }
    }
    & & \\ \\

    \theutterance \stepcounter{utterance}  
    & & & \multicolumn{2}{p{0.3\linewidth}}{
        \cellcolor[rgb]{0.9,0.9,0.9}{
            \makecell[{{p{\linewidth}}}]{
                \texttt{\tiny{[GM$|$GM]}}
                \texttt{Invalid format: tail is not empty: '} \\
\\ 
\texttt{Mijn objecten:} \\
\texttt{{-} S op (1,4)} \\
\texttt{{-} O op (4,1) } \\
\texttt{{-} T op (3,2)} \\
\texttt{{-} A op (3,5)} \\
\texttt{{-} P op (5,1)} \\
\texttt{{-} M op (6,1)} \\
\texttt{{-} D op (3,7)} \\
\\ 
\texttt{Kun jij mij jouw posities sturen? Dan kunnen we een gemeenschappelijke configuratie afspreken!} \\
            }
        }
    }
    & & \\ \\

    \theutterance \stepcounter{utterance}  
    & & \multicolumn{4}{p{0.6\linewidth}}{
        \cellcolor[rgb]{0.9,0.9,0.9}{
            \makecell[{{p{\linewidth}}}]{
                \texttt{\tiny{[P1$\langle$GM]}}
                \texttt{Straf: Je bericht mag niets na het commando bevatten!} \\
\\ 
\texttt{Zorg ervoor dat je antwoord alleen `ZEG: <BERICHT>` of `BEWEEG: <OBJECT>, (<X>, <Y>)` bevat, en niets anders!} \\
\\ 
\texttt{Jullie hebben gezamenlijk 1 van 16 straffen verzameld.} \\
\texttt{Probeer het opnieuw!} \\
            }
        }
    }
    & \\ \\

    \theutterance \stepcounter{utterance}  
    & \multicolumn{4}{p{0.6\linewidth}}{
        \cellcolor[rgb]{0.9,0.9,0.9}{
            \makecell[{{p{\linewidth}}}]{
                \texttt{\tiny{[P1$\rangle$GM]}}
                \texttt{ZEG: Hallo medespeler! Laten we onze posities vergelijken. Mijn objecten: S op (1,4), O op (4,1), T op (3,2), A op (3,5), P op (5,1), M op (6,1), D op (3,7). Wat zijn jouw posities?} \\
            }
        }
    }
    & & \\ \\

    \theutterance \stepcounter{utterance}  
    & & & \multicolumn{2}{p{0.3\linewidth}}{
        \cellcolor[rgb]{0.9,0.9,0.9}{
            \makecell[{{p{\linewidth}}}]{
                \texttt{\tiny{[GM$|$GM]}}
                \texttt{ZEG: Hallo medespeler! Laten we onze posities vergelijken. Mijn objecten: S op (1,4), O op (4,1), T op (3,2), A op (3,5), P op (5,1), M op (6,1), D op (3,7). Wat zijn jouw posities?} \\
            }
        }
    }
    & & \\ \\

    \theutterance \stepcounter{utterance}  
    & & \multicolumn{4}{p{0.6\linewidth}}{
        \cellcolor[rgb]{0.9,0.9,0.9}{
            \makecell[{{p{\linewidth}}}]{
                \texttt{\tiny{[GM$\rangle$P2]}}
                \texttt{Ik ben uw spelleider, en u speelt een samenwerkingsspel met het volgende raster als speelbord:} \\
\\ 
\texttt{```} \\
\texttt{ 1234567} \\
\texttt{╔════╤══╗} \\
\texttt{║◌◌◌◌│O◌║ 1} \\
\texttt{║P◌M◌├──╢ 2} \\
\texttt{║◌S◌◌│T◌║ 3} \\
\texttt{║◌┌──┼──╢ 4} \\
\texttt{║A│◌◌│◌◌║ 5} \\
\texttt{╟─┴┬─┤◌D║ 6} \\
\texttt{║◌◌│◌│◌◌║ 7} \\
\texttt{╚══╧═╧══╝} \\
\\ 
\texttt{```} \\
\\ 
\texttt{* De bovenrand toont x{-}coördinaten die naar rechts oplopen, en de rechterrand y{-}coördinaten die naar beneden oplopen.} \\
\texttt{* De volgende objecten zijn willekeurig op uw raster geplaatst: 'P', 'O', 'T', 'S', 'D', 'A', 'M'.} \\
\\ 
\\ 
\texttt{De andere speler ziet een variatie van het speelbord, waarbij de objecten op verschillende willekeurige locaties zijn geplaatst. U kunt het bord van de andere speler niet zien en zij kunnen het uwe niet zien.} \\
\\ 
\\ 
\texttt{**Doel:**} \\
\\ 
\texttt{Beide spelers moeten de objecten op hun respectievelijke achtergrond zo verplaatsen dat identieke objecten op dezelfde coördinaten terechtkomen. U moet communiceren met de andere speler om een gemeenschappelijke doelconfiguratie overeen te komen.} \\
\\ 
\\ 
\texttt{**Regels:**} \\
\\ 
\texttt{* In elke beurt kunt u precies één van de volgende twee commando's sturen:} \\
\texttt{ 1. `ZEG: <BERICHT>`: om een bericht (alles tot de volgende regeleinde) naar de andere speler te sturen. Ik zal het doorsturen naar uw partner.} \\
\texttt{2. `BEWEEG: <OBJECT>, (<X>, <Y>)`: om een object naar een nieuwe positie te verplaatsen, waarbij `<X>` de kolom is en `<Y>` de rij. Ik zal u laten weten of uw zet succesvol was of niet.} \\
\texttt{* Als u zich niet aan het formaat houdt, of meerdere commando's tegelijk verstuurt, moet ik u straffen.} \\
\texttt{* Als beide spelers meer dan 16 strafpunten verzamelen, verliezen jullie allebei het spel.} \\
\texttt{* Het is van vitaal belang dat u met de andere speler communiceert over uw doelstatus! De *enige* manier waarop u uw strategie aan de andere speler kunt overbrengen is door gebruik te maken van het `ZEG: <BERICHT>` commando!} \\
\\ 
\\ 
\texttt{**Objecten verplaatsen**} \\
\\ 
\texttt{* U kunt alleen objecten verplaatsen naar cellen binnen de grenzen van het raster. De doelcel moet leeg zijn, d.w.z. deze mag alleen het symbool '◌' bevatten.} \\
\texttt{* Als u probeert een object te verplaatsen naar een plek die niet leeg is, of als u probeert het buiten het raster te verplaatsen, moet ik u straffen. U krijgt een nieuwe poging.} \\
\texttt{* Voordat u een zet doet, moet u controleren of de doelplek leeg is, en geen letter, kader of lijn bevat!} \\
\\ 
\\ 
\texttt{**Einde van het spel**} \\
\\ 
\texttt{Als u denkt dat u het doel bereikt hebt om alle objecten uit te lijnen, kunt u de andere speler vragen het spel te beëindigen door `ZEG: klaar?` te sturen. Als de andere speler u vraagt om het spel te beëindigen en u antwoordt met `ZEG: klaar!`, zal het spel eindigen.} \\
\\ 
\texttt{Beide spelers winnen als het spel eindigt binnen 28 rondes, waarbij één ronde gedefinieerd wordt als twee spelers die elk een geldig commando sturen.} \\
\\ 
\texttt{**Score:**} \\
\\ 
\texttt{Hoe dichter de identieke objecten bij elkaar liggen op beide speelborden, hoe meer punten u krijgt. Strafpunten verminderen uw punten. Kunt u het record verslaan?} \\
\\ 
\texttt{De andere speler begon het spel door dit bericht te sturen:} \\
\\ 
\texttt{"Hallo medespeler! Laten we onze posities vergelijken. Mijn objecten: S op (1,4), O op (4,1), T op (3,2), A op (3,5), P op (5,1), M op (6,1), D op (3,7). Wat zijn jouw posities?"} \\
\\ 
\texttt{Wat is uw eerste opdracht?} \\
            }
        }
    }
    & \\ \\

    \theutterance \stepcounter{utterance}  
    & & & \multicolumn{4}{p{0.6\linewidth}}{
        \cellcolor[rgb]{0.9,0.9,0.9}{
            \makecell[{{p{\linewidth}}}]{
                \texttt{\tiny{[GM$\langle$P2]}}
                \texttt{Mijn eerste opdracht is:} \\
\\ 
\texttt{```} \\
\texttt{ZEG: Hallo! Mijn objecten: P op (1,2), O op (6,1), T op (6,3), S op (3,3), A op (1,5), M op (3,2), D op (6,6). Laten we een gemeenschappelijke doelconfiguratie afspreken.} \\
\texttt{```} \\
            }
        }
    }
     \\ \\

    \theutterance \stepcounter{utterance}  
    & & & \multicolumn{2}{p{0.3\linewidth}}{
        \cellcolor[rgb]{0.9,0.9,0.9}{
            \makecell[{{p{\linewidth}}}]{
                \texttt{\tiny{[GM$|$GM]}}
                \texttt{Mijn eerste opdracht is:} \\
\\ 
\texttt{```} \\
\texttt{ZEG: Hallo! Mijn objecten: P op (1,2), O op (6,1), T op (6,3), S op (3,3), A op (1,5), M op (3,2), D op (6,6). Laten we een gemeenschappelijke doelconfiguratie afspreken.} \\
\texttt{```} \\
            }
        }
    }
    & & \\ \\

    \theutterance \stepcounter{utterance}  
    & & & \multicolumn{2}{p{0.3\linewidth}}{
        \cellcolor[rgb]{0.9,0.9,0.9}{
            \makecell[{{p{\linewidth}}}]{
                \texttt{\tiny{[GM$|$GM]}}
                \texttt{Invalid format: head is not empty: 'Mijn eerste opdracht is:} \\
\\ 
\\ 
\texttt{'} \\
            }
        }
    }
    & & \\ \\

    \theutterance \stepcounter{utterance}  
    & & \multicolumn{4}{p{0.6\linewidth}}{
        \cellcolor[rgb]{0.9,0.9,0.9}{
            \makecell[{{p{\linewidth}}}]{
                \texttt{\tiny{[GM$\rangle$P2]}}
                \texttt{Straf: Je bericht mag niets vóór het commando bevatten!} \\
\\ 
\texttt{Zorg ervoor dat je antwoord alleen `ZEG: <BERICHT>` of `BEWEEG: <OBJECT>, (<X>, <Y>)` bevat, en niets anders!} \\
\\ 
\texttt{Jullie hebben gezamenlijk 2 van 16 straffen verzameld.} \\
\texttt{Probeer het opnieuw!} \\
            }
        }
    }
    & \\ \\

    \theutterance \stepcounter{utterance}  
    & & & \multicolumn{4}{p{0.6\linewidth}}{
        \cellcolor[rgb]{0.9,0.9,0.9}{
            \makecell[{{p{\linewidth}}}]{
                \texttt{\tiny{[GM$\langle$P2]}}
                \texttt{ZEG: Hallo! Mijn objecten: P op (1,2), O op (6,1), T op (6,3), S op (3,3), A op (1,5), M op (3,2), D op (6,6). Laten we een gemeenschappelijke doelconfiguratie afspreken.} \\
            }
        }
    }
     \\ \\

    \theutterance \stepcounter{utterance}  
    & & & \multicolumn{2}{p{0.3\linewidth}}{
        \cellcolor[rgb]{0.9,0.9,0.9}{
            \makecell[{{p{\linewidth}}}]{
                \texttt{\tiny{[GM$|$GM]}}
                \texttt{ZEG: Hallo! Mijn objecten: P op (1,2), O op (6,1), T op (6,3), S op (3,3), A op (1,5), M op (3,2), D op (6,6). Laten we een gemeenschappelijke doelconfiguratie afspreken.} \\
            }
        }
    }
    & & \\ \\

    \theutterance \stepcounter{utterance}  
    & & \multicolumn{4}{p{0.6\linewidth}}{
        \cellcolor[rgb]{0.9,0.9,0.9}{
            \makecell[{{p{\linewidth}}}]{
                \texttt{\tiny{[P1$\langle$GM]}}
                \texttt{Je bericht is doorgestuurd naar de andere speler.} \\
\\ 
\\ 
\texttt{Je speelt momenteel ronde 1 van maximaal 28.} \\
\texttt{Jullie hebben gezamenlijk 2 van 16 straffen verzameld.} \\
\texttt{De andere speler stuurde dit bericht:} \\
\\ 
\\ 
\texttt{"Hallo! Mijn objecten: P op (1,2), O op (6,1), T op (6,3), S op (3,3), A op (1,5), M op (3,2), D op (6,6). Laten we een gemeenschappelijke doelconfiguratie afspreken."} \\
\\ 
\\ 
\texttt{Wat is je volgende opdracht?} \\
            }
        }
    }
    & \\ \\

    \theutterance \stepcounter{utterance}  
    & \multicolumn{4}{p{0.6\linewidth}}{
        \cellcolor[rgb]{0.9,0.9,0.9}{
            \makecell[{{p{\linewidth}}}]{
                \texttt{\tiny{[P1$\rangle$GM]}}
                \texttt{ZEG: Oké, laten we een gemeenschappelijke configuratie kiezen. We moeten voor elk object dezelfde positie afspreken. Hier is een voorstel: P op (2,2), O op (2,3), T op (3,3), S op (2,4), A op (3,4), M op (4,4), D op (4,5). Wat vind jij? Of heb jij een ander voorstel?} \\
            }
        }
    }
    & & \\ \\

    \theutterance \stepcounter{utterance}  
    & & & \multicolumn{2}{p{0.3\linewidth}}{
        \cellcolor[rgb]{0.9,0.9,0.9}{
            \makecell[{{p{\linewidth}}}]{
                \texttt{\tiny{[GM$|$GM]}}
                \texttt{ZEG: Oké, laten we een gemeenschappelijke configuratie kiezen. We moeten voor elk object dezelfde positie afspreken. Hier is een voorstel: P op (2,2), O op (2,3), T op (3,3), S op (2,4), A op (3,4), M op (4,4), D op (4,5). Wat vind jij? Of heb jij een ander voorstel?} \\
            }
        }
    }
    & & \\ \\

    \theutterance \stepcounter{utterance}  
    & & \multicolumn{4}{p{0.6\linewidth}}{
        \cellcolor[rgb]{0.9,0.9,0.9}{
            \makecell[{{p{\linewidth}}}]{
                \texttt{\tiny{[GM$\rangle$P2]}}
                \texttt{Je bericht is doorgestuurd naar de andere speler.} \\
\\ 
\\ 
\texttt{Je speelt momenteel ronde 2 van maximaal 28.} \\
\texttt{Jullie hebben gezamenlijk 2 van 16 straffen verzameld.} \\
\texttt{De andere speler stuurde dit bericht:} \\
\\ 
\\ 
\texttt{"Oké, laten we een gemeenschappelijke configuratie kiezen. We moeten voor elk object dezelfde positie afspreken. Hier is een voorstel: P op (2,2), O op (2,3), T op (3,3), S op (2,4), A op (3,4), M op (4,4), D op (4,5). Wat vind jij? Of heb jij een ander voorstel?"} \\
\\ 
\\ 
\texttt{Wat is je volgende opdracht?} \\
            }
        }
    }
    & \\ \\

    \theutterance \stepcounter{utterance}  
    & & & \multicolumn{4}{p{0.6\linewidth}}{
        \cellcolor[rgb]{0.9,0.9,0.9}{
            \makecell[{{p{\linewidth}}}]{
                \texttt{\tiny{[GM$\langle$P2]}}
                \texttt{ZEG: Ik kan je voorstel niet helemaal plaatsen. Mijn bord heeft muren en lijnen, dus niet alle posities zijn bereikbaar. Ik stel voor dat we posities kiezen die op beide borden leeg en bereikbaar zijn. Laten we beginnen met één object. Is op jouw bord positie (2,3) leeg? Op mijn bord is dat een 'M'. Zullen we M op (3,2) zetten? Die is bij mij al vrij.} \\
            }
        }
    }
     \\ \\

    \theutterance \stepcounter{utterance}  
    & & & \multicolumn{2}{p{0.3\linewidth}}{
        \cellcolor[rgb]{0.9,0.9,0.9}{
            \makecell[{{p{\linewidth}}}]{
                \texttt{\tiny{[GM$|$GM]}}
                \texttt{ZEG: Ik kan je voorstel niet helemaal plaatsen. Mijn bord heeft muren en lijnen, dus niet alle posities zijn bereikbaar. Ik stel voor dat we posities kiezen die op beide borden leeg en bereikbaar zijn. Laten we beginnen met één object. Is op jouw bord positie (2,3) leeg? Op mijn bord is dat een 'M'. Zullen we M op (3,2) zetten? Die is bij mij al vrij.} \\
            }
        }
    }
    & & \\ \\

    \theutterance \stepcounter{utterance}  
    & & \multicolumn{4}{p{0.6\linewidth}}{
        \cellcolor[rgb]{0.9,0.9,0.9}{
            \makecell[{{p{\linewidth}}}]{
                \texttt{\tiny{[P1$\langle$GM]}}
                \texttt{Je bericht is doorgestuurd naar de andere speler.} \\
\\ 
\\ 
\texttt{Je speelt momenteel ronde 2 van maximaal 28.} \\
\texttt{Jullie hebben gezamenlijk 2 van 16 straffen verzameld.} \\
\texttt{De andere speler stuurde dit bericht:} \\
\\ 
\\ 
\texttt{"Ik kan je voorstel niet helemaal plaatsen. Mijn bord heeft muren en lijnen, dus niet alle posities zijn bereikbaar. Ik stel voor dat we posities kiezen die op beide borden leeg en bereikbaar zijn. Laten we beginnen met één object. Is op jouw bord positie (2,3) leeg? Op mijn bord is dat een 'M'. Zullen we M op (3,2) zetten? Die is bij mij al vrij."} \\
\\ 
\\ 
\texttt{Wat is je volgende opdracht?} \\
            }
        }
    }
    & \\ \\

    \theutterance \stepcounter{utterance}  
    & \multicolumn{4}{p{0.6\linewidth}}{
        \cellcolor[rgb]{0.9,0.9,0.9}{
            \makecell[{{p{\linewidth}}}]{
                \texttt{\tiny{[P1$\rangle$GM]}}
                \texttt{ZEG: Op mijn bord is (2,3) leeg (◌). Mijn M staat momenteel op (6,1). Jij stelt voor M op (3,2) maar op mijn bord staat daar T. Laten we posities kiezen die op beide borden leeg zijn. Is (2,2) bij jou leeg? Ik stel voor M op (2,2). Of laten we een voor een alle lege posities vergelijken. Welke posities zijn op jouw bord leeg?} \\
            }
        }
    }
    & & \\ \\

    \theutterance \stepcounter{utterance}  
    & & & \multicolumn{2}{p{0.3\linewidth}}{
        \cellcolor[rgb]{0.9,0.9,0.9}{
            \makecell[{{p{\linewidth}}}]{
                \texttt{\tiny{[GM$|$GM]}}
                \texttt{ZEG: Op mijn bord is (2,3) leeg (◌). Mijn M staat momenteel op (6,1). Jij stelt voor M op (3,2) maar op mijn bord staat daar T. Laten we posities kiezen die op beide borden leeg zijn. Is (2,2) bij jou leeg? Ik stel voor M op (2,2). Of laten we een voor een alle lege posities vergelijken. Welke posities zijn op jouw bord leeg?} \\
            }
        }
    }
    & & \\ \\

    \theutterance \stepcounter{utterance}  
    & & \multicolumn{4}{p{0.6\linewidth}}{
        \cellcolor[rgb]{0.9,0.9,0.9}{
            \makecell[{{p{\linewidth}}}]{
                \texttt{\tiny{[GM$\rangle$P2]}}
                \texttt{Je bericht is doorgestuurd naar de andere speler.} \\
\\ 
\\ 
\texttt{Je speelt momenteel ronde 3 van maximaal 28.} \\
\texttt{Jullie hebben gezamenlijk 2 van 16 straffen verzameld.} \\
\texttt{De andere speler stuurde dit bericht:} \\
\\ 
\\ 
\texttt{"Op mijn bord is (2,3) leeg (◌). Mijn M staat momenteel op (6,1). Jij stelt voor M op (3,2) maar op mijn bord staat daar T. Laten we posities kiezen die op beide borden leeg zijn. Is (2,2) bij jou leeg? Ik stel voor M op (2,2). Of laten we een voor een alle lege posities vergelijken. Welke posities zijn op jouw bord leeg?"} \\
\\ 
\\ 
\texttt{Wat is je volgende opdracht?} \\
            }
        }
    }
    & \\ \\

    \theutterance \stepcounter{utterance}  
    & & & \multicolumn{4}{p{0.6\linewidth}}{
        \cellcolor[rgb]{0.9,0.9,0.9}{
            \makecell[{{p{\linewidth}}}]{
                \texttt{\tiny{[GM$\langle$P2]}}
                \texttt{ZEG: (2,2) is bij mij leeg. M op (2,2) is prima! Andere lege posities bij mij: (2,1), (3,1), (4,1), (5,1), (7,1), (7,2), (1,3), (2,3 is M), (4,3), (5,3), (7,3), (2,4), (3,4), (4,4), (5,4), (6,4), (7,4), (2,5), (4,5), (5,5), (6,5), (7,5), (1,6), (2,6), (3,6), (4,6), (5,6), (7,6), (1,7), (2,7), (4,7), (5,7), (6,7), (7,7). Welke zijn bij jou leeg?} \\
            }
        }
    }
     \\ \\

    \theutterance \stepcounter{utterance}  
    & & & \multicolumn{2}{p{0.3\linewidth}}{
        \cellcolor[rgb]{0.9,0.9,0.9}{
            \makecell[{{p{\linewidth}}}]{
                \texttt{\tiny{[GM$|$GM]}}
                \texttt{ZEG: (2,2) is bij mij leeg. M op (2,2) is prima! Andere lege posities bij mij: (2,1), (3,1), (4,1), (5,1), (7,1), (7,2), (1,3), (2,3 is M), (4,3), (5,3), (7,3), (2,4), (3,4), (4,4), (5,4), (6,4), (7,4), (2,5), (4,5), (5,5), (6,5), (7,5), (1,6), (2,6), (3,6), (4,6), (5,6), (7,6), (1,7), (2,7), (4,7), (5,7), (6,7), (7,7). Welke zijn bij jou leeg?} \\
            }
        }
    }
    & & \\ \\

    \theutterance \stepcounter{utterance}  
    & & \multicolumn{4}{p{0.6\linewidth}}{
        \cellcolor[rgb]{0.9,0.9,0.9}{
            \makecell[{{p{\linewidth}}}]{
                \texttt{\tiny{[P1$\langle$GM]}}
                \texttt{Je bericht is doorgestuurd naar de andere speler.} \\
\\ 
\\ 
\texttt{Je speelt momenteel ronde 3 van maximaal 28.} \\
\texttt{Jullie hebben gezamenlijk 2 van 16 straffen verzameld.} \\
\texttt{De andere speler stuurde dit bericht:} \\
\\ 
\\ 
\texttt{"(2,2) is bij mij leeg. M op (2,2) is prima! Andere lege posities bij mij: (2,1), (3,1), (4,1), (5,1), (7,1), (7,2), (1,3), (2,3 is M), (4,3), (5,3), (7,3), (2,4), (3,4), (4,4), (5,4), (6,4), (7,4), (2,5), (4,5), (5,5), (6,5), (7,5), (1,6), (2,6), (3,6), (4,6), (5,6), (7,6), (1,7), (2,7), (4,7), (5,7), (6,7), (7,7). Welke zijn bij jou leeg?"} \\
\\ 
\\ 
\texttt{Wat is je volgende opdracht?} \\
            }
        }
    }
    & \\ \\

    \theutterance \stepcounter{utterance}  
    & \multicolumn{4}{p{0.6\linewidth}}{
        \cellcolor[rgb]{0.9,0.9,0.9}{
            \makecell[{{p{\linewidth}}}]{
                \texttt{\tiny{[P1$\rangle$GM]}}
                \texttt{ZEG: Mijn lege posities: (1,1), (2,1), (3,1), (1,2), (2,2), (4,2), (5,2), (6,2), (7,2), (1,3), (2,3), (3,3), (4,3), (5,3), (6,3), (7,3), (2,4), (3,4), (4,4), (5,4), (6,4), (7,4), (1,5), (2,5), (4,5), (5,5), (6,5), (7,5), (1,6), (2,6), (3,6), (4,6), (5,6), (6,6), (7,6), (1,7), (2,7), (4,7), (5,7), (6,7), (7,7). Gemeenschappelijke lege posities lijken veel. Laten we object voor object afspreken. M op (2,2) is afgesproken. Volgende: P. Waar staat jouw P nu en waar zou die heen kunnen?} \\
            }
        }
    }
    & & \\ \\

    \theutterance \stepcounter{utterance}  
    & & & \multicolumn{2}{p{0.3\linewidth}}{
        \cellcolor[rgb]{0.9,0.9,0.9}{
            \makecell[{{p{\linewidth}}}]{
                \texttt{\tiny{[GM$|$GM]}}
                \texttt{ZEG: Mijn lege posities: (1,1), (2,1), (3,1), (1,2), (2,2), (4,2), (5,2), (6,2), (7,2), (1,3), (2,3), (3,3), (4,3), (5,3), (6,3), (7,3), (2,4), (3,4), (4,4), (5,4), (6,4), (7,4), (1,5), (2,5), (4,5), (5,5), (6,5), (7,5), (1,6), (2,6), (3,6), (4,6), (5,6), (6,6), (7,6), (1,7), (2,7), (4,7), (5,7), (6,7), (7,7). Gemeenschappelijke lege posities lijken veel. Laten we object voor object afspreken. M op (2,2) is afgesproken. Volgende: P. Waar staat jouw P nu en waar zou die heen kunnen?} \\
            }
        }
    }
    & & \\ \\

    \theutterance \stepcounter{utterance}  
    & & \multicolumn{4}{p{0.6\linewidth}}{
        \cellcolor[rgb]{0.9,0.9,0.9}{
            \makecell[{{p{\linewidth}}}]{
                \texttt{\tiny{[GM$\rangle$P2]}}
                \texttt{Je bericht is doorgestuurd naar de andere speler.} \\
\\ 
\\ 
\texttt{Je speelt momenteel ronde 4 van maximaal 28.} \\
\texttt{Jullie hebben gezamenlijk 2 van 16 straffen verzameld.} \\
\texttt{De andere speler stuurde dit bericht:} \\
\\ 
\\ 
\texttt{"Mijn lege posities: (1,1), (2,1), (3,1), (1,2), (2,2), (4,2), (5,2), (6,2), (7,2), (1,3), (2,3), (3,3), (4,3), (5,3), (6,3), (7,3), (2,4), (3,4), (4,4), (5,4), (6,4), (7,4), (1,5), (2,5), (4,5), (5,5), (6,5), (7,5), (1,6), (2,6), (3,6), (4,6), (5,6), (6,6), (7,6), (1,7), (2,7), (4,7), (5,7), (6,7), (7,7). Gemeenschappelijke lege posities lijken veel. Laten we object voor object afspreken. M op (2,2) is afgesproken. Volgende: P. Waar staat jouw P nu en waar zou die heen kunnen?"} \\
\\ 
\\ 
\texttt{Wat is je volgende opdracht?} \\
            }
        }
    }
    & \\ \\

    \theutterance \stepcounter{utterance}  
    & & & \multicolumn{4}{p{0.6\linewidth}}{
        \cellcolor[rgb]{0.9,0.9,0.9}{
            \makecell[{{p{\linewidth}}}]{
                \texttt{\tiny{[GM$\langle$P2]}}
                \texttt{ZEG: M op (2,2) afgesproken! Mijn P staat op (1,2). Ik kan P verplaatsen naar bijvoorbeeld (3,1) of (4,1). Wat dacht je van P op (3,1)? Is die bij jou leeg? Of heb jij een ander voorstel voor P?} \\
            }
        }
    }
     \\ \\

    \theutterance \stepcounter{utterance}  
    & & & \multicolumn{2}{p{0.3\linewidth}}{
        \cellcolor[rgb]{0.9,0.9,0.9}{
            \makecell[{{p{\linewidth}}}]{
                \texttt{\tiny{[GM$|$GM]}}
                \texttt{ZEG: M op (2,2) afgesproken! Mijn P staat op (1,2). Ik kan P verplaatsen naar bijvoorbeeld (3,1) of (4,1). Wat dacht je van P op (3,1)? Is die bij jou leeg? Of heb jij een ander voorstel voor P?} \\
            }
        }
    }
    & & \\ \\

    \theutterance \stepcounter{utterance}  
    & & \multicolumn{4}{p{0.6\linewidth}}{
        \cellcolor[rgb]{0.9,0.9,0.9}{
            \makecell[{{p{\linewidth}}}]{
                \texttt{\tiny{[P1$\langle$GM]}}
                \texttt{Je bericht is doorgestuurd naar de andere speler.} \\
\\ 
\\ 
\texttt{Je speelt momenteel ronde 4 van maximaal 28.} \\
\texttt{Jullie hebben gezamenlijk 2 van 16 straffen verzameld.} \\
\texttt{De andere speler stuurde dit bericht:} \\
\\ 
\\ 
\texttt{"M op (2,2) afgesproken! Mijn P staat op (1,2). Ik kan P verplaatsen naar bijvoorbeeld (3,1) of (4,1). Wat dacht je van P op (3,1)? Is die bij jou leeg? Of heb jij een ander voorstel voor P?"} \\
\\ 
\\ 
\texttt{Wat is je volgende opdracht?} \\
            }
        }
    }
    & \\ \\

    \theutterance \stepcounter{utterance}  
    & \multicolumn{4}{p{0.6\linewidth}}{
        \cellcolor[rgb]{0.9,0.9,0.9}{
            \makecell[{{p{\linewidth}}}]{
                \texttt{\tiny{[P1$\rangle$GM]}}
                \texttt{ZEG: (3,1) is bij mij leeg. P op (3,1) is prima! Mijn P staat nu op (5,1). Afgesproken: M op (2,2) en P op (3,1). Volgende: O. Mijn O staat op (4,1). Waar staat jouw O en waar kunnen we die plaatsen?} \\
            }
        }
    }
    & & \\ \\

    \theutterance \stepcounter{utterance}  
    & & & \multicolumn{2}{p{0.3\linewidth}}{
        \cellcolor[rgb]{0.9,0.9,0.9}{
            \makecell[{{p{\linewidth}}}]{
                \texttt{\tiny{[GM$|$GM]}}
                \texttt{ZEG: (3,1) is bij mij leeg. P op (3,1) is prima! Mijn P staat nu op (5,1). Afgesproken: M op (2,2) en P op (3,1). Volgende: O. Mijn O staat op (4,1). Waar staat jouw O en waar kunnen we die plaatsen?} \\
            }
        }
    }
    & & \\ \\

    \theutterance \stepcounter{utterance}  
    & & \multicolumn{4}{p{0.6\linewidth}}{
        \cellcolor[rgb]{0.9,0.9,0.9}{
            \makecell[{{p{\linewidth}}}]{
                \texttt{\tiny{[GM$\rangle$P2]}}
                \texttt{Je bericht is doorgestuurd naar de andere speler.} \\
\\ 
\\ 
\texttt{Je speelt momenteel ronde 5 van maximaal 28.} \\
\texttt{Jullie hebben gezamenlijk 2 van 16 straffen verzameld.} \\
\texttt{De andere speler stuurde dit bericht:} \\
\\ 
\\ 
\texttt{"(3,1) is bij mij leeg. P op (3,1) is prima! Mijn P staat nu op (5,1). Afgesproken: M op (2,2) en P op (3,1). Volgende: O. Mijn O staat op (4,1). Waar staat jouw O en waar kunnen we die plaatsen?"} \\
\\ 
\\ 
\texttt{Wat is je volgende opdracht?} \\
            }
        }
    }
    & \\ \\

    \theutterance \stepcounter{utterance}  
    & & & \multicolumn{4}{p{0.6\linewidth}}{
        \cellcolor[rgb]{0.9,0.9,0.9}{
            \makecell[{{p{\linewidth}}}]{
                \texttt{\tiny{[GM$\langle$P2]}}
                \texttt{ZEG: Mijn O staat op (6,1). O op (4,1) is bij mij leeg, dus dat lijkt me prima! Afgesproken: M op (2,2), P op (3,1), O op (4,1). Volgende: T. Mijn T staat op (6,3). Wat dacht je van T op (3,3)? Is die bij jou leeg? Of heb jij een ander voorstel?} \\
            }
        }
    }
     \\ \\

    \theutterance \stepcounter{utterance}  
    & & & \multicolumn{2}{p{0.3\linewidth}}{
        \cellcolor[rgb]{0.9,0.9,0.9}{
            \makecell[{{p{\linewidth}}}]{
                \texttt{\tiny{[GM$|$GM]}}
                \texttt{ZEG: Mijn O staat op (6,1). O op (4,1) is bij mij leeg, dus dat lijkt me prima! Afgesproken: M op (2,2), P op (3,1), O op (4,1). Volgende: T. Mijn T staat op (6,3). Wat dacht je van T op (3,3)? Is die bij jou leeg? Of heb jij een ander voorstel?} \\
            }
        }
    }
    & & \\ \\

    \theutterance \stepcounter{utterance}  
    & & \multicolumn{4}{p{0.6\linewidth}}{
        \cellcolor[rgb]{0.9,0.9,0.9}{
            \makecell[{{p{\linewidth}}}]{
                \texttt{\tiny{[P1$\langle$GM]}}
                \texttt{Je bericht is doorgestuurd naar de andere speler.} \\
\\ 
\\ 
\texttt{Je speelt momenteel ronde 5 van maximaal 28.} \\
\texttt{Jullie hebben gezamenlijk 2 van 16 straffen verzameld.} \\
\texttt{De andere speler stuurde dit bericht:} \\
\\ 
\\ 
\texttt{"Mijn O staat op (6,1). O op (4,1) is bij mij leeg, dus dat lijkt me prima! Afgesproken: M op (2,2), P op (3,1), O op (4,1). Volgende: T. Mijn T staat op (6,3). Wat dacht je van T op (3,3)? Is die bij jou leeg? Of heb jij een ander voorstel?"} \\
\\ 
\\ 
\texttt{Wat is je volgende opdracht?} \\
            }
        }
    }
    & \\ \\

    \theutterance \stepcounter{utterance}  
    & \multicolumn{4}{p{0.6\linewidth}}{
        \cellcolor[rgb]{0.9,0.9,0.9}{
            \makecell[{{p{\linewidth}}}]{
                \texttt{\tiny{[P1$\rangle$GM]}}
                \texttt{ZEG: (3,3) is bij mij leeg. T op (3,3) is goed! Mijn T staat nu op (3,2). Afgesproken: M(2,2), P(3,1), O(4,1), T(3,3). Volgende: S. Mijn S staat op (1,4). Waar staat jouw S?} \\
            }
        }
    }
    & & \\ \\

    \theutterance \stepcounter{utterance}  
    & & & \multicolumn{2}{p{0.3\linewidth}}{
        \cellcolor[rgb]{0.9,0.9,0.9}{
            \makecell[{{p{\linewidth}}}]{
                \texttt{\tiny{[GM$|$GM]}}
                \texttt{ZEG: (3,3) is bij mij leeg. T op (3,3) is goed! Mijn T staat nu op (3,2). Afgesproken: M(2,2), P(3,1), O(4,1), T(3,3). Volgende: S. Mijn S staat op (1,4). Waar staat jouw S?} \\
            }
        }
    }
    & & \\ \\

    \theutterance \stepcounter{utterance}  
    & & \multicolumn{4}{p{0.6\linewidth}}{
        \cellcolor[rgb]{0.9,0.9,0.9}{
            \makecell[{{p{\linewidth}}}]{
                \texttt{\tiny{[GM$\rangle$P2]}}
                \texttt{Je bericht is doorgestuurd naar de andere speler.} \\
\\ 
\\ 
\texttt{Je speelt momenteel ronde 6 van maximaal 28.} \\
\texttt{Jullie hebben gezamenlijk 2 van 16 straffen verzameld.} \\
\texttt{De andere speler stuurde dit bericht:} \\
\\ 
\\ 
\texttt{"(3,3) is bij mij leeg. T op (3,3) is goed! Mijn T staat nu op (3,2). Afgesproken: M(2,2), P(3,1), O(4,1), T(3,3). Volgende: S. Mijn S staat op (1,4). Waar staat jouw S?"} \\
\\ 
\\ 
\texttt{Wat is je volgende opdracht?} \\
            }
        }
    }
    & \\ \\

    \theutterance \stepcounter{utterance}  
    & & & \multicolumn{4}{p{0.6\linewidth}}{
        \cellcolor[rgb]{0.9,0.9,0.9}{
            \makecell[{{p{\linewidth}}}]{
                \texttt{\tiny{[GM$\langle$P2]}}
                \texttt{ZEG: Mijn S staat op (3,3), maar dat is nu T. Wat dacht je van S op (2,4)? Is die bij jou leeg? Of heb jij een ander voorstel voor S?} \\
            }
        }
    }
     \\ \\

    \theutterance \stepcounter{utterance}  
    & & & \multicolumn{2}{p{0.3\linewidth}}{
        \cellcolor[rgb]{0.9,0.9,0.9}{
            \makecell[{{p{\linewidth}}}]{
                \texttt{\tiny{[GM$|$GM]}}
                \texttt{ZEG: Mijn S staat op (3,3), maar dat is nu T. Wat dacht je van S op (2,4)? Is die bij jou leeg? Of heb jij een ander voorstel voor S?} \\
            }
        }
    }
    & & \\ \\

    \theutterance \stepcounter{utterance}  
    & & \multicolumn{4}{p{0.6\linewidth}}{
        \cellcolor[rgb]{0.9,0.9,0.9}{
            \makecell[{{p{\linewidth}}}]{
                \texttt{\tiny{[P1$\langle$GM]}}
                \texttt{Je bericht is doorgestuurd naar de andere speler.} \\
\\ 
\\ 
\texttt{Je speelt momenteel ronde 6 van maximaal 28.} \\
\texttt{Jullie hebben gezamenlijk 2 van 16 straffen verzameld.} \\
\texttt{De andere speler stuurde dit bericht:} \\
\\ 
\\ 
\texttt{"Mijn S staat op (3,3), maar dat is nu T. Wat dacht je van S op (2,4)? Is die bij jou leeg? Of heb jij een ander voorstel voor S?"} \\
\\ 
\\ 
\texttt{Wat is je volgende opdracht?} \\
            }
        }
    }
    & \\ \\

    \theutterance \stepcounter{utterance}  
    & \multicolumn{4}{p{0.6\linewidth}}{
        \cellcolor[rgb]{0.9,0.9,0.9}{
            \makecell[{{p{\linewidth}}}]{
                \texttt{\tiny{[P1$\rangle$GM]}}
                \texttt{ZEG: (2,4) is bij mij leeg. S op (2,4) is prima! Afgesproken: M(2,2), P(3,1), O(4,1), T(3,3), S(2,4). Volgende: A. Mijn A staat op (3,5). Waar staat jouw A?} \\
            }
        }
    }
    & & \\ \\

    \theutterance \stepcounter{utterance}  
    & & & \multicolumn{2}{p{0.3\linewidth}}{
        \cellcolor[rgb]{0.9,0.9,0.9}{
            \makecell[{{p{\linewidth}}}]{
                \texttt{\tiny{[GM$|$GM]}}
                \texttt{ZEG: (2,4) is bij mij leeg. S op (2,4) is prima! Afgesproken: M(2,2), P(3,1), O(4,1), T(3,3), S(2,4). Volgende: A. Mijn A staat op (3,5). Waar staat jouw A?} \\
            }
        }
    }
    & & \\ \\

    \theutterance \stepcounter{utterance}  
    & & \multicolumn{4}{p{0.6\linewidth}}{
        \cellcolor[rgb]{0.9,0.9,0.9}{
            \makecell[{{p{\linewidth}}}]{
                \texttt{\tiny{[GM$\rangle$P2]}}
                \texttt{Je bericht is doorgestuurd naar de andere speler.} \\
\\ 
\\ 
\texttt{Je speelt momenteel ronde 7 van maximaal 28.} \\
\texttt{Jullie hebben gezamenlijk 2 van 16 straffen verzameld.} \\
\texttt{De andere speler stuurde dit bericht:} \\
\\ 
\\ 
\texttt{"(2,4) is bij mij leeg. S op (2,4) is prima! Afgesproken: M(2,2), P(3,1), O(4,1), T(3,3), S(2,4). Volgende: A. Mijn A staat op (3,5). Waar staat jouw A?"} \\
\\ 
\\ 
\texttt{Wat is je volgende opdracht?} \\
            }
        }
    }
    & \\ \\

    \theutterance \stepcounter{utterance}  
    & & & \multicolumn{4}{p{0.6\linewidth}}{
        \cellcolor[rgb]{0.9,0.9,0.9}{
            \makecell[{{p{\linewidth}}}]{
                \texttt{\tiny{[GM$\langle$P2]}}
                \texttt{ZEG: Mijn A staat op (1,5). Is (2,5) bij jou leeg? Anders kunnen we A ook op (3,4) zetten. Wat denk jij?} \\
            }
        }
    }
     \\ \\

    \theutterance \stepcounter{utterance}  
    & & & \multicolumn{2}{p{0.3\linewidth}}{
        \cellcolor[rgb]{0.9,0.9,0.9}{
            \makecell[{{p{\linewidth}}}]{
                \texttt{\tiny{[GM$|$GM]}}
                \texttt{ZEG: Mijn A staat op (1,5). Is (2,5) bij jou leeg? Anders kunnen we A ook op (3,4) zetten. Wat denk jij?} \\
            }
        }
    }
    & & \\ \\

    \theutterance \stepcounter{utterance}  
    & & \multicolumn{4}{p{0.6\linewidth}}{
        \cellcolor[rgb]{0.9,0.9,0.9}{
            \makecell[{{p{\linewidth}}}]{
                \texttt{\tiny{[P1$\langle$GM]}}
                \texttt{Je bericht is doorgestuurd naar de andere speler.} \\
\\ 
\\ 
\texttt{Je speelt momenteel ronde 7 van maximaal 28.} \\
\texttt{Jullie hebben gezamenlijk 2 van 16 straffen verzameld.} \\
\texttt{De andere speler stuurde dit bericht:} \\
\\ 
\\ 
\texttt{"Mijn A staat op (1,5). Is (2,5) bij jou leeg? Anders kunnen we A ook op (3,4) zetten. Wat denk jij?"} \\
\\ 
\\ 
\texttt{Wat is je volgende opdracht?} \\
            }
        }
    }
    & \\ \\

    \theutterance \stepcounter{utterance}  
    & \multicolumn{4}{p{0.6\linewidth}}{
        \cellcolor[rgb]{0.9,0.9,0.9}{
            \makecell[{{p{\linewidth}}}]{
                \texttt{\tiny{[P1$\rangle$GM]}}
                \texttt{ZEG: Zowel (2,5) als (3,4) zijn bij mij leeg. Laten we A op (3,4) zetten, dan staan de letters wat dichter bij elkaar. Afgesproken: M(2,2), P(3,1), O(4,1), T(3,3), S(2,4), A(3,4). Nu nog D. Mijn D staat op (3,7). Waar staat jouw D?} \\
            }
        }
    }
    & & \\ \\

    \theutterance \stepcounter{utterance}  
    & & & \multicolumn{2}{p{0.3\linewidth}}{
        \cellcolor[rgb]{0.9,0.9,0.9}{
            \makecell[{{p{\linewidth}}}]{
                \texttt{\tiny{[GM$|$GM]}}
                \texttt{ZEG: Zowel (2,5) als (3,4) zijn bij mij leeg. Laten we A op (3,4) zetten, dan staan de letters wat dichter bij elkaar. Afgesproken: M(2,2), P(3,1), O(4,1), T(3,3), S(2,4), A(3,4). Nu nog D. Mijn D staat op (3,7). Waar staat jouw D?} \\
            }
        }
    }
    & & \\ \\

    \theutterance \stepcounter{utterance}  
    & & \multicolumn{4}{p{0.6\linewidth}}{
        \cellcolor[rgb]{0.9,0.9,0.9}{
            \makecell[{{p{\linewidth}}}]{
                \texttt{\tiny{[GM$\rangle$P2]}}
                \texttt{Je bericht is doorgestuurd naar de andere speler.} \\
\\ 
\\ 
\texttt{Je speelt momenteel ronde 8 van maximaal 28.} \\
\texttt{Jullie hebben gezamenlijk 2 van 16 straffen verzameld.} \\
\texttt{De andere speler stuurde dit bericht:} \\
\\ 
\\ 
\texttt{"Zowel (2,5) als (3,4) zijn bij mij leeg. Laten we A op (3,4) zetten, dan staan de letters wat dichter bij elkaar. Afgesproken: M(2,2), P(3,1), O(4,1), T(3,3), S(2,4), A(3,4). Nu nog D. Mijn D staat op (3,7). Waar staat jouw D?"} \\
\\ 
\\ 
\texttt{Wat is je volgende opdracht?} \\
            }
        }
    }
    & \\ \\

    \theutterance \stepcounter{utterance}  
    & & & \multicolumn{4}{p{0.6\linewidth}}{
        \cellcolor[rgb]{0.9,0.9,0.9}{
            \makecell[{{p{\linewidth}}}]{
                \texttt{\tiny{[GM$\langle$P2]}}
                \texttt{ZEG: Mijn D staat op (6,6). Is (4,5) bij jou leeg? Dan kunnen we D daar zetten. Afgesproken: M(2,2), P(3,1), O(4,1), T(3,3), S(2,4), A(3,4), D(4,5). Als je akkoord bent, kunnen we beginnen met verplaatsen!} \\
            }
        }
    }
     \\ \\

    \theutterance \stepcounter{utterance}  
    & & & \multicolumn{2}{p{0.3\linewidth}}{
        \cellcolor[rgb]{0.9,0.9,0.9}{
            \makecell[{{p{\linewidth}}}]{
                \texttt{\tiny{[GM$|$GM]}}
                \texttt{ZEG: Mijn D staat op (6,6). Is (4,5) bij jou leeg? Dan kunnen we D daar zetten. Afgesproken: M(2,2), P(3,1), O(4,1), T(3,3), S(2,4), A(3,4), D(4,5). Als je akkoord bent, kunnen we beginnen met verplaatsen!} \\
            }
        }
    }
    & & \\ \\

    \theutterance \stepcounter{utterance}  
    & & \multicolumn{4}{p{0.6\linewidth}}{
        \cellcolor[rgb]{0.9,0.9,0.9}{
            \makecell[{{p{\linewidth}}}]{
                \texttt{\tiny{[P1$\langle$GM]}}
                \texttt{Je bericht is doorgestuurd naar de andere speler.} \\
\\ 
\\ 
\texttt{Je speelt momenteel ronde 8 van maximaal 28.} \\
\texttt{Jullie hebben gezamenlijk 2 van 16 straffen verzameld.} \\
\texttt{De andere speler stuurde dit bericht:} \\
\\ 
\\ 
\texttt{"Mijn D staat op (6,6). Is (4,5) bij jou leeg? Dan kunnen we D daar zetten. Afgesproken: M(2,2), P(3,1), O(4,1), T(3,3), S(2,4), A(3,4), D(4,5). Als je akkoord bent, kunnen we beginnen met verplaatsen!"} \\
\\ 
\\ 
\texttt{Wat is je volgende opdracht?} \\
            }
        }
    }
    & \\ \\

    \theutterance \stepcounter{utterance}  
    & \multicolumn{4}{p{0.6\linewidth}}{
        \cellcolor[rgb]{0.9,0.9,0.9}{
            \makecell[{{p{\linewidth}}}]{
                \texttt{\tiny{[P1$\rangle$GM]}}
                \texttt{ZEG: (4,5) is bij mij leeg. Helemaal akkoord met de configuratie! Laten we beginnen met verplaatsen. Ik zet eerst M naar (2,2). Succes!} \\
            }
        }
    }
    & & \\ \\

    \theutterance \stepcounter{utterance}  
    & & & \multicolumn{2}{p{0.3\linewidth}}{
        \cellcolor[rgb]{0.9,0.9,0.9}{
            \makecell[{{p{\linewidth}}}]{
                \texttt{\tiny{[GM$|$GM]}}
                \texttt{ZEG: (4,5) is bij mij leeg. Helemaal akkoord met de configuratie! Laten we beginnen met verplaatsen. Ik zet eerst M naar (2,2). Succes!} \\
            }
        }
    }
    & & \\ \\

    \theutterance \stepcounter{utterance}  
    & & \multicolumn{4}{p{0.6\linewidth}}{
        \cellcolor[rgb]{0.9,0.9,0.9}{
            \makecell[{{p{\linewidth}}}]{
                \texttt{\tiny{[GM$\rangle$P2]}}
                \texttt{Je bericht is doorgestuurd naar de andere speler.} \\
\\ 
\\ 
\texttt{Je speelt momenteel ronde 9 van maximaal 28.} \\
\texttt{Jullie hebben gezamenlijk 2 van 16 straffen verzameld.} \\
\texttt{De andere speler stuurde dit bericht:} \\
\\ 
\\ 
\texttt{"(4,5) is bij mij leeg. Helemaal akkoord met de configuratie! Laten we beginnen met verplaatsen. Ik zet eerst M naar (2,2). Succes!"} \\
\\ 
\\ 
\texttt{Wat is je volgende opdracht?} \\
            }
        }
    }
    & \\ \\

    \theutterance \stepcounter{utterance}  
    & & & \multicolumn{4}{p{0.6\linewidth}}{
        \cellcolor[rgb]{0.9,0.9,0.9}{
            \makecell[{{p{\linewidth}}}]{
                \texttt{\tiny{[GM$\langle$P2]}}
                \texttt{BEWEEG: M, (2,2)} \\
            }
        }
    }
     \\ \\

    \theutterance \stepcounter{utterance}  
    & & & \multicolumn{2}{p{0.3\linewidth}}{
        \cellcolor[rgb]{0.9,0.9,0.9}{
            \makecell[{{p{\linewidth}}}]{
                \texttt{\tiny{[GM$|$GM]}}
                \texttt{BEWEEG: M, (2,2)} \\
            }
        }
    }
    & & \\ \\

    \theutterance \stepcounter{utterance}  
    & & & \multicolumn{2}{p{0.3\linewidth}}{
        \cellcolor[rgb]{0.9,0.9,0.9}{
            \makecell[{{p{\linewidth}}}]{
                \texttt{\tiny{[GM$|$GM]}}
                \texttt{M moved to (2, 2)} \\
            }
        }
    }
    & & \\ \\

    \theutterance \stepcounter{utterance}  
    & & \multicolumn{4}{p{0.6\linewidth}}{
        \cellcolor[rgb]{0.9,0.9,0.9}{
            \makecell[{{p{\linewidth}}}]{
                \texttt{\tiny{[P1$\langle$GM]}}
                \texttt{Je bericht is doorgestuurd naar de andere speler.} \\
\\ 
\\ 
\texttt{Je speelt momenteel ronde 9 van maximaal 28.} \\
\texttt{Jullie hebben gezamenlijk 2 van 16 straffen verzameld.} \\
\texttt{De andere speler heeft een voorwerp op zijn rooster verplaatst.} \\
\\ 
\\ 
\texttt{Wat is je volgende opdracht?} \\
            }
        }
    }
    & \\ \\

    \theutterance \stepcounter{utterance}  
    & \multicolumn{4}{p{0.6\linewidth}}{
        \cellcolor[rgb]{0.9,0.9,0.9}{
            \makecell[{{p{\linewidth}}}]{
                \texttt{\tiny{[P1$\rangle$GM]}}
                \texttt{BEWEEG: M, (2,2)} \\
            }
        }
    }
    & & \\ \\

    \theutterance \stepcounter{utterance}  
    & & & \multicolumn{2}{p{0.3\linewidth}}{
        \cellcolor[rgb]{0.9,0.9,0.9}{
            \makecell[{{p{\linewidth}}}]{
                \texttt{\tiny{[GM$|$GM]}}
                \texttt{BEWEEG: M, (2,2)} \\
            }
        }
    }
    & & \\ \\

    \theutterance \stepcounter{utterance}  
    & & & \multicolumn{2}{p{0.3\linewidth}}{
        \cellcolor[rgb]{0.9,0.9,0.9}{
            \makecell[{{p{\linewidth}}}]{
                \texttt{\tiny{[GM$|$GM]}}
                \texttt{M moved to (2, 2)} \\
            }
        }
    }
    & & \\ \\

    \theutterance \stepcounter{utterance}  
    & & \multicolumn{4}{p{0.6\linewidth}}{
        \cellcolor[rgb]{0.9,0.9,0.9}{
            \makecell[{{p{\linewidth}}}]{
                \texttt{\tiny{[GM$\rangle$P2]}}
                \texttt{Verplaatst 'M' naar (2,2) met succes. Uw bijgewerkte raster ziet er als volgt uit:} \\
\texttt{```} \\
\texttt{ 1234567} \\
\texttt{╔════╤══╗} \\
\texttt{║◌◌◌◌│O◌║ 1} \\
\texttt{║PM◌◌├──╢ 2} \\
\texttt{║◌S◌◌│T◌║ 3} \\
\texttt{║◌┌──┼──╢ 4} \\
\texttt{║A│◌◌│◌◌║ 5} \\
\texttt{╟─┴┬─┤◌D║ 6} \\
\texttt{║◌◌│◌│◌◌║ 7} \\
\texttt{╚══╧═╧══╝} \\
\texttt{```} \\
\texttt{Je speelt momenteel ronde 10 van maximaal 28.} \\
\texttt{Jullie hebben gezamenlijk 2 van 16 straffen verzameld.} \\
\texttt{De andere speler heeft een voorwerp op zijn rooster verplaatst.} \\
\\ 
\\ 
\texttt{Wat is je volgende opdracht?} \\
            }
        }
    }
    & \\ \\

    \theutterance \stepcounter{utterance}  
    & & & \multicolumn{4}{p{0.6\linewidth}}{
        \cellcolor[rgb]{0.9,0.9,0.9}{
            \makecell[{{p{\linewidth}}}]{
                \texttt{\tiny{[GM$\langle$P2]}}
                \texttt{BEWEEG: P, (3,1)} \\
            }
        }
    }
     \\ \\

    \theutterance \stepcounter{utterance}  
    & & & \multicolumn{2}{p{0.3\linewidth}}{
        \cellcolor[rgb]{0.9,0.9,0.9}{
            \makecell[{{p{\linewidth}}}]{
                \texttt{\tiny{[GM$|$GM]}}
                \texttt{BEWEEG: P, (3,1)} \\
            }
        }
    }
    & & \\ \\

    \theutterance \stepcounter{utterance}  
    & & & \multicolumn{2}{p{0.3\linewidth}}{
        \cellcolor[rgb]{0.9,0.9,0.9}{
            \makecell[{{p{\linewidth}}}]{
                \texttt{\tiny{[GM$|$GM]}}
                \texttt{P moved to (3, 1)} \\
            }
        }
    }
    & & \\ \\

    \theutterance \stepcounter{utterance}  
    & & \multicolumn{4}{p{0.6\linewidth}}{
        \cellcolor[rgb]{0.9,0.9,0.9}{
            \makecell[{{p{\linewidth}}}]{
                \texttt{\tiny{[P1$\langle$GM]}}
                \texttt{Verplaatst 'M' naar (2,2) met succes. Uw bijgewerkte raster ziet er als volgt uit:} \\
\texttt{```} \\
\texttt{ 1234567} \\
\texttt{╔════╤══╗} \\
\texttt{║◌◌◌O│P◌║ 1} \\
\texttt{║◌MT◌├──╢ 2} \\
\texttt{║◌◌◌◌│◌◌║ 3} \\
\texttt{║S┌──┼──╢ 4} \\
\texttt{║◌│◌A│◌◌║ 5} \\
\texttt{╟─┴┬─┤◌◌║ 6} \\
\texttt{║◌◌│◌│D◌║ 7} \\
\texttt{╚══╧═╧══╝} \\
\texttt{```} \\
\texttt{Je speelt momenteel ronde 10 van maximaal 28.} \\
\texttt{Jullie hebben gezamenlijk 2 van 16 straffen verzameld.} \\
\texttt{De andere speler heeft een voorwerp op zijn rooster verplaatst.} \\
\\ 
\\ 
\texttt{Wat is je volgende opdracht?} \\
            }
        }
    }
    & \\ \\

    \theutterance \stepcounter{utterance}  
    & \multicolumn{4}{p{0.6\linewidth}}{
        \cellcolor[rgb]{0.9,0.9,0.9}{
            \makecell[{{p{\linewidth}}}]{
                \texttt{\tiny{[P1$\rangle$GM]}}
                \texttt{BEWEEG: P, (3,1)} \\
            }
        }
    }
    & & \\ \\

    \theutterance \stepcounter{utterance}  
    & & & \multicolumn{2}{p{0.3\linewidth}}{
        \cellcolor[rgb]{0.9,0.9,0.9}{
            \makecell[{{p{\linewidth}}}]{
                \texttt{\tiny{[GM$|$GM]}}
                \texttt{BEWEEG: P, (3,1)} \\
            }
        }
    }
    & & \\ \\

    \theutterance \stepcounter{utterance}  
    & & & \multicolumn{2}{p{0.3\linewidth}}{
        \cellcolor[rgb]{0.9,0.9,0.9}{
            \makecell[{{p{\linewidth}}}]{
                \texttt{\tiny{[GM$|$GM]}}
                \texttt{P moved to (3, 1)} \\
            }
        }
    }
    & & \\ \\

    \theutterance \stepcounter{utterance}  
    & & \multicolumn{4}{p{0.6\linewidth}}{
        \cellcolor[rgb]{0.9,0.9,0.9}{
            \makecell[{{p{\linewidth}}}]{
                \texttt{\tiny{[GM$\rangle$P2]}}
                \texttt{Verplaatst 'P' naar (3,1) met succes. Uw bijgewerkte raster ziet er als volgt uit:} \\
\texttt{```} \\
\texttt{ 1234567} \\
\texttt{╔════╤══╗} \\
\texttt{║◌◌P◌│O◌║ 1} \\
\texttt{║◌M◌◌├──╢ 2} \\
\texttt{║◌S◌◌│T◌║ 3} \\
\texttt{║◌┌──┼──╢ 4} \\
\texttt{║A│◌◌│◌◌║ 5} \\
\texttt{╟─┴┬─┤◌D║ 6} \\
\texttt{║◌◌│◌│◌◌║ 7} \\
\texttt{╚══╧═╧══╝} \\
\texttt{```} \\
\texttt{Je speelt momenteel ronde 11 van maximaal 28.} \\
\texttt{Jullie hebben gezamenlijk 2 van 16 straffen verzameld.} \\
\texttt{De andere speler heeft een voorwerp op zijn rooster verplaatst.} \\
\\ 
\\ 
\texttt{Wat is je volgende opdracht?} \\
            }
        }
    }
    & \\ \\

    \theutterance \stepcounter{utterance}  
    & & & \multicolumn{4}{p{0.6\linewidth}}{
        \cellcolor[rgb]{0.9,0.9,0.9}{
            \makecell[{{p{\linewidth}}}]{
                \texttt{\tiny{[GM$\langle$P2]}}
                \texttt{BEWEEG: O, (4,1)} \\
            }
        }
    }
     \\ \\

    \theutterance \stepcounter{utterance}  
    & & & \multicolumn{2}{p{0.3\linewidth}}{
        \cellcolor[rgb]{0.9,0.9,0.9}{
            \makecell[{{p{\linewidth}}}]{
                \texttt{\tiny{[GM$|$GM]}}
                \texttt{BEWEEG: O, (4,1)} \\
            }
        }
    }
    & & \\ \\

    \theutterance \stepcounter{utterance}  
    & & & \multicolumn{2}{p{0.3\linewidth}}{
        \cellcolor[rgb]{0.9,0.9,0.9}{
            \makecell[{{p{\linewidth}}}]{
                \texttt{\tiny{[GM$|$GM]}}
                \texttt{O moved to (4, 1)} \\
            }
        }
    }
    & & \\ \\

    \theutterance \stepcounter{utterance}  
    & & \multicolumn{4}{p{0.6\linewidth}}{
        \cellcolor[rgb]{0.9,0.9,0.9}{
            \makecell[{{p{\linewidth}}}]{
                \texttt{\tiny{[P1$\langle$GM]}}
                \texttt{Verplaatst 'P' naar (3,1) met succes. Uw bijgewerkte raster ziet er als volgt uit:} \\
\texttt{```} \\
\texttt{ 1234567} \\
\texttt{╔════╤══╗} \\
\texttt{║◌◌PO│◌◌║ 1} \\
\texttt{║◌MT◌├──╢ 2} \\
\texttt{║◌◌◌◌│◌◌║ 3} \\
\texttt{║S┌──┼──╢ 4} \\
\texttt{║◌│◌A│◌◌║ 5} \\
\texttt{╟─┴┬─┤◌◌║ 6} \\
\texttt{║◌◌│◌│D◌║ 7} \\
\texttt{╚══╧═╧══╝} \\
\texttt{```} \\
\texttt{Je speelt momenteel ronde 11 van maximaal 28.} \\
\texttt{Jullie hebben gezamenlijk 2 van 16 straffen verzameld.} \\
\texttt{De andere speler heeft een voorwerp op zijn rooster verplaatst.} \\
\\ 
\\ 
\texttt{Wat is je volgende opdracht?} \\
            }
        }
    }
    & \\ \\

    \theutterance \stepcounter{utterance}  
    & \multicolumn{4}{p{0.6\linewidth}}{
        \cellcolor[rgb]{0.9,0.9,0.9}{
            \makecell[{{p{\linewidth}}}]{
                \texttt{\tiny{[P1$\rangle$GM]}}
                \texttt{BEWEEG: O, (4,1)} \\
            }
        }
    }
    & & \\ \\

    \theutterance \stepcounter{utterance}  
    & & & \multicolumn{2}{p{0.3\linewidth}}{
        \cellcolor[rgb]{0.9,0.9,0.9}{
            \makecell[{{p{\linewidth}}}]{
                \texttt{\tiny{[GM$|$GM]}}
                \texttt{BEWEEG: O, (4,1)} \\
            }
        }
    }
    & & \\ \\

    \theutterance \stepcounter{utterance}  
    & & & \multicolumn{2}{p{0.3\linewidth}}{
        \cellcolor[rgb]{0.9,0.9,0.9}{
            \makecell[{{p{\linewidth}}}]{
                \texttt{\tiny{[GM$|$GM]}}
                \texttt{Straf: (4,1) is niet leeg, maar bevat 'O'.} \\
\\ 
\texttt{Jullie hebben gezamenlijk 3 van 16 straffen verzameld.} \\
\texttt{Probeer het opnieuw!} \\
            }
        }
    }
    & & \\ \\

    \theutterance \stepcounter{utterance}  
    & & \multicolumn{4}{p{0.6\linewidth}}{
        \cellcolor[rgb]{0.9,0.9,0.9}{
            \makecell[{{p{\linewidth}}}]{
                \texttt{\tiny{[P1$\langle$GM]}}
                \texttt{Straf: (4,1) is niet leeg, maar bevat 'O'.} \\
\\ 
\texttt{Jullie hebben gezamenlijk 3 van 16 straffen verzameld.} \\
\texttt{Probeer het opnieuw!} \\
            }
        }
    }
    & \\ \\

    \theutterance \stepcounter{utterance}  
    & \multicolumn{4}{p{0.6\linewidth}}{
        \cellcolor[rgb]{0.9,0.9,0.9}{
            \makecell[{{p{\linewidth}}}]{
                \texttt{\tiny{[P1$\rangle$GM]}}
                \texttt{BEWEEG: T, (3,3)} \\
            }
        }
    }
    & & \\ \\

    \theutterance \stepcounter{utterance}  
    & & & \multicolumn{2}{p{0.3\linewidth}}{
        \cellcolor[rgb]{0.9,0.9,0.9}{
            \makecell[{{p{\linewidth}}}]{
                \texttt{\tiny{[GM$|$GM]}}
                \texttt{BEWEEG: T, (3,3)} \\
            }
        }
    }
    & & \\ \\

    \theutterance \stepcounter{utterance}  
    & & & \multicolumn{2}{p{0.3\linewidth}}{
        \cellcolor[rgb]{0.9,0.9,0.9}{
            \makecell[{{p{\linewidth}}}]{
                \texttt{\tiny{[GM$|$GM]}}
                \texttt{T moved to (3, 3)} \\
            }
        }
    }
    & & \\ \\

    \theutterance \stepcounter{utterance}  
    & & \multicolumn{4}{p{0.6\linewidth}}{
        \cellcolor[rgb]{0.9,0.9,0.9}{
            \makecell[{{p{\linewidth}}}]{
                \texttt{\tiny{[GM$\rangle$P2]}}
                \texttt{Verplaatst 'O' naar (4,1) met succes. Uw bijgewerkte raster ziet er als volgt uit:} \\
\texttt{```} \\
\texttt{ 1234567} \\
\texttt{╔════╤══╗} \\
\texttt{║◌◌PO│◌◌║ 1} \\
\texttt{║◌M◌◌├──╢ 2} \\
\texttt{║◌S◌◌│T◌║ 3} \\
\texttt{║◌┌──┼──╢ 4} \\
\texttt{║A│◌◌│◌◌║ 5} \\
\texttt{╟─┴┬─┤◌D║ 6} \\
\texttt{║◌◌│◌│◌◌║ 7} \\
\texttt{╚══╧═╧══╝} \\
\texttt{```} \\
\texttt{Je speelt momenteel ronde 12 van maximaal 28.} \\
\texttt{Jullie hebben gezamenlijk 3 van 16 straffen verzameld.} \\
\texttt{De andere speler heeft een voorwerp op zijn rooster verplaatst.} \\
\\ 
\\ 
\texttt{Wat is je volgende opdracht?} \\
            }
        }
    }
    & \\ \\

    \theutterance \stepcounter{utterance}  
    & & & \multicolumn{4}{p{0.6\linewidth}}{
        \cellcolor[rgb]{0.9,0.9,0.9}{
            \makecell[{{p{\linewidth}}}]{
                \texttt{\tiny{[GM$\langle$P2]}}
                \texttt{BEWEEG: T, (3,3)} \\
            }
        }
    }
     \\ \\

    \theutterance \stepcounter{utterance}  
    & & & \multicolumn{2}{p{0.3\linewidth}}{
        \cellcolor[rgb]{0.9,0.9,0.9}{
            \makecell[{{p{\linewidth}}}]{
                \texttt{\tiny{[GM$|$GM]}}
                \texttt{BEWEEG: T, (3,3)} \\
            }
        }
    }
    & & \\ \\

    \theutterance \stepcounter{utterance}  
    & & & \multicolumn{2}{p{0.3\linewidth}}{
        \cellcolor[rgb]{0.9,0.9,0.9}{
            \makecell[{{p{\linewidth}}}]{
                \texttt{\tiny{[GM$|$GM]}}
                \texttt{T moved to (3, 3)} \\
            }
        }
    }
    & & \\ \\

    \theutterance \stepcounter{utterance}  
    & & \multicolumn{4}{p{0.6\linewidth}}{
        \cellcolor[rgb]{0.9,0.9,0.9}{
            \makecell[{{p{\linewidth}}}]{
                \texttt{\tiny{[P1$\langle$GM]}}
                \texttt{Verplaatst 'T' naar (3,3) met succes. Uw bijgewerkte raster ziet er als volgt uit:} \\
\texttt{```} \\
\texttt{ 1234567} \\
\texttt{╔════╤══╗} \\
\texttt{║◌◌PO│◌◌║ 1} \\
\texttt{║◌M◌◌├──╢ 2} \\
\texttt{║◌◌T◌│◌◌║ 3} \\
\texttt{║S┌──┼──╢ 4} \\
\texttt{║◌│◌A│◌◌║ 5} \\
\texttt{╟─┴┬─┤◌◌║ 6} \\
\texttt{║◌◌│◌│D◌║ 7} \\
\texttt{╚══╧═╧══╝} \\
\texttt{```} \\
\texttt{Je speelt momenteel ronde 12 van maximaal 28.} \\
\texttt{Jullie hebben gezamenlijk 3 van 16 straffen verzameld.} \\
\texttt{De andere speler heeft een voorwerp op zijn rooster verplaatst.} \\
\\ 
\\ 
\texttt{Wat is je volgende opdracht?} \\
            }
        }
    }
    & \\ \\

    \theutterance \stepcounter{utterance}  
    & \multicolumn{4}{p{0.6\linewidth}}{
        \cellcolor[rgb]{0.9,0.9,0.9}{
            \makecell[{{p{\linewidth}}}]{
                \texttt{\tiny{[P1$\rangle$GM]}}
                \texttt{BEWEEG: S, (2,4)} \\
            }
        }
    }
    & & \\ \\

    \theutterance \stepcounter{utterance}  
    & & & \multicolumn{2}{p{0.3\linewidth}}{
        \cellcolor[rgb]{0.9,0.9,0.9}{
            \makecell[{{p{\linewidth}}}]{
                \texttt{\tiny{[GM$|$GM]}}
                \texttt{BEWEEG: S, (2,4)} \\
            }
        }
    }
    & & \\ \\

    \theutterance \stepcounter{utterance}  
    & & & \multicolumn{2}{p{0.3\linewidth}}{
        \cellcolor[rgb]{0.9,0.9,0.9}{
            \makecell[{{p{\linewidth}}}]{
                \texttt{\tiny{[GM$|$GM]}}
                \texttt{Straf: (2,4) is niet leeg, maar bevat '┌'.} \\
\\ 
\texttt{Jullie hebben gezamenlijk 4 van 16 straffen verzameld.} \\
\texttt{Probeer het opnieuw!} \\
            }
        }
    }
    & & \\ \\

    \theutterance \stepcounter{utterance}  
    & & \multicolumn{4}{p{0.6\linewidth}}{
        \cellcolor[rgb]{0.9,0.9,0.9}{
            \makecell[{{p{\linewidth}}}]{
                \texttt{\tiny{[P1$\langle$GM]}}
                \texttt{Straf: (2,4) is niet leeg, maar bevat '┌'.} \\
\\ 
\texttt{Jullie hebben gezamenlijk 4 van 16 straffen verzameld.} \\
\texttt{Probeer het opnieuw!} \\
            }
        }
    }
    & \\ \\

    \theutterance \stepcounter{utterance}  
    & \multicolumn{4}{p{0.6\linewidth}}{
        \cellcolor[rgb]{0.9,0.9,0.9}{
            \makecell[{{p{\linewidth}}}]{
                \texttt{\tiny{[P1$\rangle$GM]}}
                \texttt{BEWEEG: S, (2,4)} \\
            }
        }
    }
    & & \\ \\

    \theutterance \stepcounter{utterance}  
    & & & \multicolumn{2}{p{0.3\linewidth}}{
        \cellcolor[rgb]{0.9,0.9,0.9}{
            \makecell[{{p{\linewidth}}}]{
                \texttt{\tiny{[GM$|$GM]}}
                \texttt{BEWEEG: S, (2,4)} \\
            }
        }
    }
    & & \\ \\

    \theutterance \stepcounter{utterance}  
    & & & \multicolumn{2}{p{0.3\linewidth}}{
        \cellcolor[rgb]{0.9,0.9,0.9}{
            \makecell[{{p{\linewidth}}}]{
                \texttt{\tiny{[GM$|$GM]}}
                \texttt{Straf: (2,4) is niet leeg, maar bevat '┌'.} \\
\\ 
\texttt{Jullie hebben gezamenlijk 5 van 16 straffen verzameld.} \\
\texttt{Probeer het opnieuw!} \\
            }
        }
    }
    & & \\ \\

    \theutterance \stepcounter{utterance}  
    & & \multicolumn{4}{p{0.6\linewidth}}{
        \cellcolor[rgb]{0.9,0.9,0.9}{
            \makecell[{{p{\linewidth}}}]{
                \texttt{\tiny{[P1$\langle$GM]}}
                \texttt{Straf: (2,4) is niet leeg, maar bevat '┌'.} \\
\\ 
\texttt{Jullie hebben gezamenlijk 5 van 16 straffen verzameld.} \\
\texttt{Probeer het opnieuw!} \\
            }
        }
    }
    & \\ \\

    \theutterance \stepcounter{utterance}  
    & \multicolumn{4}{p{0.6\linewidth}}{
        \cellcolor[rgb]{0.9,0.9,0.9}{
            \makecell[{{p{\linewidth}}}]{
                \texttt{\tiny{[P1$\rangle$GM]}}
                \texttt{BEWEEG: S, (3,4)} \\
            }
        }
    }
    & & \\ \\

    \theutterance \stepcounter{utterance}  
    & & & \multicolumn{2}{p{0.3\linewidth}}{
        \cellcolor[rgb]{0.9,0.9,0.9}{
            \makecell[{{p{\linewidth}}}]{
                \texttt{\tiny{[GM$|$GM]}}
                \texttt{BEWEEG: S, (3,4)} \\
            }
        }
    }
    & & \\ \\

    \theutterance \stepcounter{utterance}  
    & & & \multicolumn{2}{p{0.3\linewidth}}{
        \cellcolor[rgb]{0.9,0.9,0.9}{
            \makecell[{{p{\linewidth}}}]{
                \texttt{\tiny{[GM$|$GM]}}
                \texttt{Straf: (3,4) is niet leeg, maar bevat '─'.} \\
\\ 
\texttt{Jullie hebben gezamenlijk 6 van 16 straffen verzameld.} \\
\texttt{Probeer het opnieuw!} \\
            }
        }
    }
    & & \\ \\

    \theutterance \stepcounter{utterance}  
    & & \multicolumn{4}{p{0.6\linewidth}}{
        \cellcolor[rgb]{0.9,0.9,0.9}{
            \makecell[{{p{\linewidth}}}]{
                \texttt{\tiny{[P1$\langle$GM]}}
                \texttt{Straf: (3,4) is niet leeg, maar bevat '─'.} \\
\\ 
\texttt{Jullie hebben gezamenlijk 6 van 16 straffen verzameld.} \\
\texttt{Probeer het opnieuw!} \\
            }
        }
    }
    & \\ \\

    \theutterance \stepcounter{utterance}  
    & \multicolumn{4}{p{0.6\linewidth}}{
        \cellcolor[rgb]{0.9,0.9,0.9}{
            \makecell[{{p{\linewidth}}}]{
                \texttt{\tiny{[P1$\rangle$GM]}}
                \texttt{BEWEEG: S, (5,4)} \\
            }
        }
    }
    & & \\ \\

    \theutterance \stepcounter{utterance}  
    & & & \multicolumn{2}{p{0.3\linewidth}}{
        \cellcolor[rgb]{0.9,0.9,0.9}{
            \makecell[{{p{\linewidth}}}]{
                \texttt{\tiny{[GM$|$GM]}}
                \texttt{BEWEEG: S, (5,4)} \\
            }
        }
    }
    & & \\ \\

    \theutterance \stepcounter{utterance}  
    & & & \multicolumn{2}{p{0.3\linewidth}}{
        \cellcolor[rgb]{0.9,0.9,0.9}{
            \makecell[{{p{\linewidth}}}]{
                \texttt{\tiny{[GM$|$GM]}}
                \texttt{Straf: (5,4) is niet leeg, maar bevat '┼'.} \\
\\ 
\texttt{Jullie hebben gezamenlijk 7 van 16 straffen verzameld.} \\
\texttt{Probeer het opnieuw!} \\
            }
        }
    }
    & & \\ \\

    \theutterance \stepcounter{utterance}  
    & & \multicolumn{4}{p{0.6\linewidth}}{
        \cellcolor[rgb]{0.9,0.9,0.9}{
            \makecell[{{p{\linewidth}}}]{
                \texttt{\tiny{[P1$\langle$GM]}}
                \texttt{Straf: (5,4) is niet leeg, maar bevat '┼'.} \\
\\ 
\texttt{Jullie hebben gezamenlijk 7 van 16 straffen verzameld.} \\
\texttt{Probeer het opnieuw!} \\
            }
        }
    }
    & \\ \\

    \theutterance \stepcounter{utterance}  
    & \multicolumn{4}{p{0.6\linewidth}}{
        \cellcolor[rgb]{0.9,0.9,0.9}{
            \makecell[{{p{\linewidth}}}]{
                \texttt{\tiny{[P1$\rangle$GM]}}
                \texttt{ZEG: Let op, ik kan S niet naar (2,4) verplaatsen want daar zit een muurdeel '┌'. Kun jij me vertellen welke posities in rij 4 bij jou leeg zijn? Mijn rij 4 heeft alleen (1,4), (6,4) en (7,4) leeg. Komt (2,4) bij jou wel uit?} \\
            }
        }
    }
    & & \\ \\

    \theutterance \stepcounter{utterance}  
    & & & \multicolumn{2}{p{0.3\linewidth}}{
        \cellcolor[rgb]{0.9,0.9,0.9}{
            \makecell[{{p{\linewidth}}}]{
                \texttt{\tiny{[GM$|$GM]}}
                \texttt{ZEG: Let op, ik kan S niet naar (2,4) verplaatsen want daar zit een muurdeel '┌'. Kun jij me vertellen welke posities in rij 4 bij jou leeg zijn? Mijn rij 4 heeft alleen (1,4), (6,4) en (7,4) leeg. Komt (2,4) bij jou wel uit?} \\
            }
        }
    }
    & & \\ \\

    \theutterance \stepcounter{utterance}  
    & & \multicolumn{4}{p{0.6\linewidth}}{
        \cellcolor[rgb]{0.9,0.9,0.9}{
            \makecell[{{p{\linewidth}}}]{
                \texttt{\tiny{[GM$\rangle$P2]}}
                \texttt{Verplaatst 'T' naar (3,3) met succes. Uw bijgewerkte raster ziet er als volgt uit:} \\
\texttt{```} \\
\texttt{ 1234567} \\
\texttt{╔════╤══╗} \\
\texttt{║◌◌PO│◌◌║ 1} \\
\texttt{║◌M◌◌├──╢ 2} \\
\texttt{║◌ST◌│◌◌║ 3} \\
\texttt{║◌┌──┼──╢ 4} \\
\texttt{║A│◌◌│◌◌║ 5} \\
\texttt{╟─┴┬─┤◌D║ 6} \\
\texttt{║◌◌│◌│◌◌║ 7} \\
\texttt{╚══╧═╧══╝} \\
\texttt{```} \\
\texttt{Je speelt momenteel ronde 13 van maximaal 28.} \\
\texttt{Jullie hebben gezamenlijk 7 van 16 straffen verzameld.} \\
\texttt{De andere speler stuurde dit bericht:} \\
\\ 
\\ 
\texttt{"Let op, ik kan S niet naar (2,4) verplaatsen want daar zit een muurdeel '┌'. Kun jij me vertellen welke posities in rij 4 bij jou leeg zijn? Mijn rij 4 heeft alleen (1,4), (6,4) en (7,4) leeg. Komt (2,4) bij jou wel uit?"} \\
\\ 
\\ 
\texttt{Wat is je volgende opdracht?} \\
            }
        }
    }
    & \\ \\

    \theutterance \stepcounter{utterance}  
    & & & \multicolumn{4}{p{0.6\linewidth}}{
        \cellcolor[rgb]{0.9,0.9,0.9}{
            \makecell[{{p{\linewidth}}}]{
                \texttt{\tiny{[GM$\langle$P2]}}
                \texttt{ZEG: Bij mij is (2,4) leeg, maar bij jou dus niet. Laten we S op een andere plek zetten die voor ons allebei leeg is. Is (1,4) bij jou leeg? Mijn (1,4) heeft een muurdeel '┌', dus die is niet leeg. Wat dacht je van S op (3,4)? Maar die is afgesproken voor A. Zullen we S en A ruilen? Dus S op (3,4) en A op (2,4)? Of heb jij een ander voorstel?} \\
            }
        }
    }
     \\ \\

    \theutterance \stepcounter{utterance}  
    & & & \multicolumn{2}{p{0.3\linewidth}}{
        \cellcolor[rgb]{0.9,0.9,0.9}{
            \makecell[{{p{\linewidth}}}]{
                \texttt{\tiny{[GM$|$GM]}}
                \texttt{ZEG: Bij mij is (2,4) leeg, maar bij jou dus niet. Laten we S op een andere plek zetten die voor ons allebei leeg is. Is (1,4) bij jou leeg? Mijn (1,4) heeft een muurdeel '┌', dus die is niet leeg. Wat dacht je van S op (3,4)? Maar die is afgesproken voor A. Zullen we S en A ruilen? Dus S op (3,4) en A op (2,4)? Of heb jij een ander voorstel?} \\
            }
        }
    }
    & & \\ \\

    \theutterance \stepcounter{utterance}  
    & & \multicolumn{4}{p{0.6\linewidth}}{
        \cellcolor[rgb]{0.9,0.9,0.9}{
            \makecell[{{p{\linewidth}}}]{
                \texttt{\tiny{[P1$\langle$GM]}}
                \texttt{Je bericht is doorgestuurd naar de andere speler.} \\
\\ 
\\ 
\texttt{Je speelt momenteel ronde 13 van maximaal 28.} \\
\texttt{Jullie hebben gezamenlijk 7 van 16 straffen verzameld.} \\
\texttt{De andere speler stuurde dit bericht:} \\
\\ 
\\ 
\texttt{"Bij mij is (2,4) leeg, maar bij jou dus niet. Laten we S op een andere plek zetten die voor ons allebei leeg is. Is (1,4) bij jou leeg? Mijn (1,4) heeft een muurdeel '┌', dus die is niet leeg. Wat dacht je van S op (3,4)? Maar die is afgesproken voor A. Zullen we S en A ruilen? Dus S op (3,4) en A op (2,4)? Of heb jij een ander voorstel?"} \\
\\ 
\\ 
\texttt{Wat is je volgende opdracht?} \\
            }
        }
    }
    & \\ \\

    \theutterance \stepcounter{utterance}  
    & \multicolumn{4}{p{0.6\linewidth}}{
        \cellcolor[rgb]{0.9,0.9,0.9}{
            \makecell[{{p{\linewidth}}}]{
                \texttt{\tiny{[P1$\rangle$GM]}}
                \texttt{ZEG: Goed idee, laten we ruilen. S op (3,4) en A op (2,4). Maar wacht, bij mij is (3,4) ook een muurdeel '─'. En (2,4) is '┌'. Rij 4 is bij mij lastig. Wat zijn bij jou de lege posities in rijen 4 en 5? Bij mij zijn (1,4), (6,4), (7,4) leeg, en in rij 5: (1,5), (2,5), (4,5), (5,5), (6,5), (7,5). Zullen we nieuwe posities voor S, A en D kiezen op basis hiervan?} \\
            }
        }
    }
    & & \\ \\

    \theutterance \stepcounter{utterance}  
    & & & \multicolumn{2}{p{0.3\linewidth}}{
        \cellcolor[rgb]{0.9,0.9,0.9}{
            \makecell[{{p{\linewidth}}}]{
                \texttt{\tiny{[GM$|$GM]}}
                \texttt{ZEG: Goed idee, laten we ruilen. S op (3,4) en A op (2,4). Maar wacht, bij mij is (3,4) ook een muurdeel '─'. En (2,4) is '┌'. Rij 4 is bij mij lastig. Wat zijn bij jou de lege posities in rijen 4 en 5? Bij mij zijn (1,4), (6,4), (7,4) leeg, en in rij 5: (1,5), (2,5), (4,5), (5,5), (6,5), (7,5). Zullen we nieuwe posities voor S, A en D kiezen op basis hiervan?} \\
            }
        }
    }
    & & \\ \\

    \theutterance \stepcounter{utterance}  
    & & \multicolumn{4}{p{0.6\linewidth}}{
        \cellcolor[rgb]{0.9,0.9,0.9}{
            \makecell[{{p{\linewidth}}}]{
                \texttt{\tiny{[GM$\rangle$P2]}}
                \texttt{Je bericht is doorgestuurd naar de andere speler.} \\
\\ 
\\ 
\texttt{Je speelt momenteel ronde 14 van maximaal 28.} \\
\texttt{Jullie hebben gezamenlijk 7 van 16 straffen verzameld.} \\
\texttt{De andere speler stuurde dit bericht:} \\
\\ 
\\ 
\texttt{"Goed idee, laten we ruilen. S op (3,4) en A op (2,4). Maar wacht, bij mij is (3,4) ook een muurdeel '─'. En (2,4) is '┌'. Rij 4 is bij mij lastig. Wat zijn bij jou de lege posities in rijen 4 en 5? Bij mij zijn (1,4), (6,4), (7,4) leeg, en in rij 5: (1,5), (2,5), (4,5), (5,5), (6,5), (7,5). Zullen we nieuwe posities voor S, A en D kiezen op basis hiervan?"} \\
\\ 
\\ 
\texttt{Wat is je volgende opdracht?} \\
            }
        }
    }
    & \\ \\

    \theutterance \stepcounter{utterance}  
    & & & \multicolumn{4}{p{0.6\linewidth}}{
        \cellcolor[rgb]{0.9,0.9,0.9}{
            \makecell[{{p{\linewidth}}}]{
                \texttt{\tiny{[GM$\langle$P2]}}
                \texttt{ZEG: Bij mij leeg in rij 4: (2,4), (3,4), (4,4), (5,4), (6,4), (7,4). In rij 5 leeg: (2,5), (4,5), (5,5), (6,5), (7,5). Gemeenschappelijk leeg in rij 4: (6,4), (7,4). In rij 5: (2,5), (4,5), (5,5), (6,5), (7,5). Voorstel: S op (6,4), A op (2,5), D op (4,5) zoals eerder. Of zullen we S, A en D alledrie in rij 5 zetten? Bijv. S(2,5), A(4,5), D(5,5)?} \\
            }
        }
    }
     \\ \\

    \theutterance \stepcounter{utterance}  
    & & & \multicolumn{2}{p{0.3\linewidth}}{
        \cellcolor[rgb]{0.9,0.9,0.9}{
            \makecell[{{p{\linewidth}}}]{
                \texttt{\tiny{[GM$|$GM]}}
                \texttt{ZEG: Bij mij leeg in rij 4: (2,4), (3,4), (4,4), (5,4), (6,4), (7,4). In rij 5 leeg: (2,5), (4,5), (5,5), (6,5), (7,5). Gemeenschappelijk leeg in rij 4: (6,4), (7,4). In rij 5: (2,5), (4,5), (5,5), (6,5), (7,5). Voorstel: S op (6,4), A op (2,5), D op (4,5) zoals eerder. Of zullen we S, A en D alledrie in rij 5 zetten? Bijv. S(2,5), A(4,5), D(5,5)?} \\
            }
        }
    }
    & & \\ \\

    \theutterance \stepcounter{utterance}  
    & & \multicolumn{4}{p{0.6\linewidth}}{
        \cellcolor[rgb]{0.9,0.9,0.9}{
            \makecell[{{p{\linewidth}}}]{
                \texttt{\tiny{[P1$\langle$GM]}}
                \texttt{Je bericht is doorgestuurd naar de andere speler.} \\
\\ 
\\ 
\texttt{Je speelt momenteel ronde 14 van maximaal 28.} \\
\texttt{Jullie hebben gezamenlijk 7 van 16 straffen verzameld.} \\
\texttt{De andere speler stuurde dit bericht:} \\
\\ 
\\ 
\texttt{"Bij mij leeg in rij 4: (2,4), (3,4), (4,4), (5,4), (6,4), (7,4). In rij 5 leeg: (2,5), (4,5), (5,5), (6,5), (7,5). Gemeenschappelijk leeg in rij 4: (6,4), (7,4). In rij 5: (2,5), (4,5), (5,5), (6,5), (7,5). Voorstel: S op (6,4), A op (2,5), D op (4,5) zoals eerder. Of zullen we S, A en D alledrie in rij 5 zetten? Bijv. S(2,5), A(4,5), D(5,5)?"} \\
\\ 
\\ 
\texttt{Wat is je volgende opdracht?} \\
            }
        }
    }
    & \\ \\

    \theutterance \stepcounter{utterance}  
    & \multicolumn{4}{p{0.6\linewidth}}{
        \cellcolor[rgb]{0.9,0.9,0.9}{
            \makecell[{{p{\linewidth}}}]{
                \texttt{\tiny{[P1$\rangle$GM]}}
                \texttt{ZEG: Laten we S(6,4), A(2,5), D(4,5) doen. Die posities zijn bij mij ook leeg. M(2,2), P(3,1), O(4,1) en T(3,3) staan al goed. Dan blijft het overzichtelijk. Akkoord?} \\
            }
        }
    }
    & & \\ \\

    \theutterance \stepcounter{utterance}  
    & & & \multicolumn{2}{p{0.3\linewidth}}{
        \cellcolor[rgb]{0.9,0.9,0.9}{
            \makecell[{{p{\linewidth}}}]{
                \texttt{\tiny{[GM$|$GM]}}
                \texttt{ZEG: Laten we S(6,4), A(2,5), D(4,5) doen. Die posities zijn bij mij ook leeg. M(2,2), P(3,1), O(4,1) en T(3,3) staan al goed. Dan blijft het overzichtelijk. Akkoord?} \\
            }
        }
    }
    & & \\ \\

    \theutterance \stepcounter{utterance}  
    & & \multicolumn{4}{p{0.6\linewidth}}{
        \cellcolor[rgb]{0.9,0.9,0.9}{
            \makecell[{{p{\linewidth}}}]{
                \texttt{\tiny{[GM$\rangle$P2]}}
                \texttt{Je bericht is doorgestuurd naar de andere speler.} \\
\\ 
\\ 
\texttt{Je speelt momenteel ronde 15 van maximaal 28.} \\
\texttt{Jullie hebben gezamenlijk 7 van 16 straffen verzameld.} \\
\texttt{De andere speler stuurde dit bericht:} \\
\\ 
\\ 
\texttt{"Laten we S(6,4), A(2,5), D(4,5) doen. Die posities zijn bij mij ook leeg. M(2,2), P(3,1), O(4,1) en T(3,3) staan al goed. Dan blijft het overzichtelijk. Akkoord?"} \\
\\ 
\\ 
\texttt{Wat is je volgende opdracht?} \\
            }
        }
    }
    & \\ \\

    \theutterance \stepcounter{utterance}  
    & & & \multicolumn{4}{p{0.6\linewidth}}{
        \cellcolor[rgb]{0.9,0.9,0.9}{
            \makecell[{{p{\linewidth}}}]{
                \texttt{\tiny{[GM$\langle$P2]}}
                \texttt{ZEG: Akkoord! Configuratie: M(2,2), P(3,1), O(4,1), T(3,3), S(6,4), A(2,5), D(4,5). Ik ga nu S verplaatsen naar (6,4).} \\
            }
        }
    }
     \\ \\

    \theutterance \stepcounter{utterance}  
    & & & \multicolumn{2}{p{0.3\linewidth}}{
        \cellcolor[rgb]{0.9,0.9,0.9}{
            \makecell[{{p{\linewidth}}}]{
                \texttt{\tiny{[GM$|$GM]}}
                \texttt{ZEG: Akkoord! Configuratie: M(2,2), P(3,1), O(4,1), T(3,3), S(6,4), A(2,5), D(4,5). Ik ga nu S verplaatsen naar (6,4).} \\
            }
        }
    }
    & & \\ \\

    \theutterance \stepcounter{utterance}  
    & & \multicolumn{4}{p{0.6\linewidth}}{
        \cellcolor[rgb]{0.9,0.9,0.9}{
            \makecell[{{p{\linewidth}}}]{
                \texttt{\tiny{[P1$\langle$GM]}}
                \texttt{Je bericht is doorgestuurd naar de andere speler.} \\
\\ 
\\ 
\texttt{Je speelt momenteel ronde 15 van maximaal 28.} \\
\texttt{Jullie hebben gezamenlijk 7 van 16 straffen verzameld.} \\
\texttt{De andere speler stuurde dit bericht:} \\
\\ 
\\ 
\texttt{"Akkoord! Configuratie: M(2,2), P(3,1), O(4,1), T(3,3), S(6,4), A(2,5), D(4,5). Ik ga nu S verplaatsen naar (6,4)."} \\
\\ 
\\ 
\texttt{Wat is je volgende opdracht?} \\
            }
        }
    }
    & \\ \\

    \theutterance \stepcounter{utterance}  
    & \multicolumn{4}{p{0.6\linewidth}}{
        \cellcolor[rgb]{0.9,0.9,0.9}{
            \makecell[{{p{\linewidth}}}]{
                \texttt{\tiny{[P1$\rangle$GM]}}
                \texttt{BEWEEG: S, (6,4)} \\
            }
        }
    }
    & & \\ \\

    \theutterance \stepcounter{utterance}  
    & & & \multicolumn{2}{p{0.3\linewidth}}{
        \cellcolor[rgb]{0.9,0.9,0.9}{
            \makecell[{{p{\linewidth}}}]{
                \texttt{\tiny{[GM$|$GM]}}
                \texttt{BEWEEG: S, (6,4)} \\
            }
        }
    }
    & & \\ \\

    \theutterance \stepcounter{utterance}  
    & & & \multicolumn{2}{p{0.3\linewidth}}{
        \cellcolor[rgb]{0.9,0.9,0.9}{
            \makecell[{{p{\linewidth}}}]{
                \texttt{\tiny{[GM$|$GM]}}
                \texttt{Straf: (6,4) is niet leeg, maar bevat '─'.} \\
\\ 
\texttt{Jullie hebben gezamenlijk 8 van 16 straffen verzameld.} \\
\texttt{Probeer het opnieuw!} \\
            }
        }
    }
    & & \\ \\

    \theutterance \stepcounter{utterance}  
    & & \multicolumn{4}{p{0.6\linewidth}}{
        \cellcolor[rgb]{0.9,0.9,0.9}{
            \makecell[{{p{\linewidth}}}]{
                \texttt{\tiny{[P1$\langle$GM]}}
                \texttt{Straf: (6,4) is niet leeg, maar bevat '─'.} \\
\\ 
\texttt{Jullie hebben gezamenlijk 8 van 16 straffen verzameld.} \\
\texttt{Probeer het opnieuw!} \\
            }
        }
    }
    & \\ \\

    \theutterance \stepcounter{utterance}  
    & \multicolumn{4}{p{0.6\linewidth}}{
        \cellcolor[rgb]{0.9,0.9,0.9}{
            \makecell[{{p{\linewidth}}}]{
                \texttt{\tiny{[P1$\rangle$GM]}}
                \texttt{BEWEEG: S, (7,4)} \\
            }
        }
    }
    & & \\ \\

    \theutterance \stepcounter{utterance}  
    & & & \multicolumn{2}{p{0.3\linewidth}}{
        \cellcolor[rgb]{0.9,0.9,0.9}{
            \makecell[{{p{\linewidth}}}]{
                \texttt{\tiny{[GM$|$GM]}}
                \texttt{BEWEEG: S, (7,4)} \\
            }
        }
    }
    & & \\ \\

    \theutterance \stepcounter{utterance}  
    & & & \multicolumn{2}{p{0.3\linewidth}}{
        \cellcolor[rgb]{0.9,0.9,0.9}{
            \makecell[{{p{\linewidth}}}]{
                \texttt{\tiny{[GM$|$GM]}}
                \texttt{Straf: (7,4) is niet leeg, maar bevat '─'.} \\
\\ 
\texttt{Jullie hebben gezamenlijk 9 van 16 straffen verzameld.} \\
\texttt{Probeer het opnieuw!} \\
            }
        }
    }
    & & \\ \\

    \theutterance \stepcounter{utterance}  
    & & \multicolumn{4}{p{0.6\linewidth}}{
        \cellcolor[rgb]{0.9,0.9,0.9}{
            \makecell[{{p{\linewidth}}}]{
                \texttt{\tiny{[P1$\langle$GM]}}
                \texttt{Straf: (7,4) is niet leeg, maar bevat '─'.} \\
\\ 
\texttt{Jullie hebben gezamenlijk 9 van 16 straffen verzameld.} \\
\texttt{Probeer het opnieuw!} \\
            }
        }
    }
    & \\ \\

    \theutterance \stepcounter{utterance}  
    & \multicolumn{4}{p{0.6\linewidth}}{
        \cellcolor[rgb]{0.9,0.9,0.9}{
            \makecell[{{p{\linewidth}}}]{
                \texttt{\tiny{[P1$\rangle$GM]}}
                \texttt{ZEG: Sorry, ik vergiste me. Bij mij zijn (6,4) en (7,4) ook muren '─'. Mijn enige lege in rij 4 is (1,4) waar S al staat. Maar jij zei dat (1,4) bij jou een muur is. Dus S kan niet in rij 4. Kunnen we S, A en D alledrie in rij 5 zetten, waar we veel lege gemeenschappelijke posities hebben? Voorstel: S(2,5), A(5,5), D(4,5) {-} allemaal leeg bij mij. Wat denk jij?} \\
            }
        }
    }
    & & \\ \\

    \theutterance \stepcounter{utterance}  
    & & & \multicolumn{2}{p{0.3\linewidth}}{
        \cellcolor[rgb]{0.9,0.9,0.9}{
            \makecell[{{p{\linewidth}}}]{
                \texttt{\tiny{[GM$|$GM]}}
                \texttt{ZEG: Sorry, ik vergiste me. Bij mij zijn (6,4) en (7,4) ook muren '─'. Mijn enige lege in rij 4 is (1,4) waar S al staat. Maar jij zei dat (1,4) bij jou een muur is. Dus S kan niet in rij 4. Kunnen we S, A en D alledrie in rij 5 zetten, waar we veel lege gemeenschappelijke posities hebben? Voorstel: S(2,5), A(5,5), D(4,5) {-} allemaal leeg bij mij. Wat denk jij?} \\
            }
        }
    }
    & & \\ \\

    \theutterance \stepcounter{utterance}  
    & & \multicolumn{4}{p{0.6\linewidth}}{
        \cellcolor[rgb]{0.9,0.9,0.9}{
            \makecell[{{p{\linewidth}}}]{
                \texttt{\tiny{[GM$\rangle$P2]}}
                \texttt{Je bericht is doorgestuurd naar de andere speler.} \\
\\ 
\\ 
\texttt{Je speelt momenteel ronde 16 van maximaal 28.} \\
\texttt{Jullie hebben gezamenlijk 9 van 16 straffen verzameld.} \\
\texttt{De andere speler stuurde dit bericht:} \\
\\ 
\\ 
\texttt{"Sorry, ik vergiste me. Bij mij zijn (6,4) en (7,4) ook muren '─'. Mijn enige lege in rij 4 is (1,4) waar S al staat. Maar jij zei dat (1,4) bij jou een muur is. Dus S kan niet in rij 4. Kunnen we S, A en D alledrie in rij 5 zetten, waar we veel lege gemeenschappelijke posities hebben? Voorstel: S(2,5), A(5,5), D(4,5) {-} allemaal leeg bij mij. Wat denk jij?"} \\
\\ 
\\ 
\texttt{Wat is je volgende opdracht?} \\
            }
        }
    }
    & \\ \\

    \theutterance \stepcounter{utterance}  
    & & & \multicolumn{4}{p{0.6\linewidth}}{
        \cellcolor[rgb]{0.9,0.9,0.9}{
            \makecell[{{p{\linewidth}}}]{
                \texttt{\tiny{[GM$\langle$P2]}}
                \texttt{ZEG: Geen probleem! S(2,5), A(5,5), D(4,5) is prima. Bij mij zijn die allemaal leeg. Nieuwe configuratie: M(2,2), P(3,1), O(4,1), T(3,3), S(2,5), A(5,5), D(4,5). Ik ga S naar (2,5) verplaatsen.} \\
            }
        }
    }
     \\ \\

    \theutterance \stepcounter{utterance}  
    & & & \multicolumn{2}{p{0.3\linewidth}}{
        \cellcolor[rgb]{0.9,0.9,0.9}{
            \makecell[{{p{\linewidth}}}]{
                \texttt{\tiny{[GM$|$GM]}}
                \texttt{ZEG: Geen probleem! S(2,5), A(5,5), D(4,5) is prima. Bij mij zijn die allemaal leeg. Nieuwe configuratie: M(2,2), P(3,1), O(4,1), T(3,3), S(2,5), A(5,5), D(4,5). Ik ga S naar (2,5) verplaatsen.} \\
            }
        }
    }
    & & \\ \\

    \theutterance \stepcounter{utterance}  
    & & \multicolumn{4}{p{0.6\linewidth}}{
        \cellcolor[rgb]{0.9,0.9,0.9}{
            \makecell[{{p{\linewidth}}}]{
                \texttt{\tiny{[P1$\langle$GM]}}
                \texttt{Je bericht is doorgestuurd naar de andere speler.} \\
\\ 
\\ 
\texttt{Je speelt momenteel ronde 16 van maximaal 28.} \\
\texttt{Jullie hebben gezamenlijk 9 van 16 straffen verzameld.} \\
\texttt{De andere speler stuurde dit bericht:} \\
\\ 
\\ 
\texttt{"Geen probleem! S(2,5), A(5,5), D(4,5) is prima. Bij mij zijn die allemaal leeg. Nieuwe configuratie: M(2,2), P(3,1), O(4,1), T(3,3), S(2,5), A(5,5), D(4,5). Ik ga S naar (2,5) verplaatsen."} \\
\\ 
\\ 
\texttt{Wat is je volgende opdracht?} \\
            }
        }
    }
    & \\ \\

    \theutterance \stepcounter{utterance}  
    & \multicolumn{4}{p{0.6\linewidth}}{
        \cellcolor[rgb]{0.9,0.9,0.9}{
            \makecell[{{p{\linewidth}}}]{
                \texttt{\tiny{[P1$\rangle$GM]}}
                \texttt{BEWEEG: S, (2,5)} \\
            }
        }
    }
    & & \\ \\

    \theutterance \stepcounter{utterance}  
    & & & \multicolumn{2}{p{0.3\linewidth}}{
        \cellcolor[rgb]{0.9,0.9,0.9}{
            \makecell[{{p{\linewidth}}}]{
                \texttt{\tiny{[GM$|$GM]}}
                \texttt{BEWEEG: S, (2,5)} \\
            }
        }
    }
    & & \\ \\

    \theutterance \stepcounter{utterance}  
    & & & \multicolumn{2}{p{0.3\linewidth}}{
        \cellcolor[rgb]{0.9,0.9,0.9}{
            \makecell[{{p{\linewidth}}}]{
                \texttt{\tiny{[GM$|$GM]}}
                \texttt{Straf: (2,5) is niet leeg, maar bevat '│'.} \\
\\ 
\texttt{Jullie hebben gezamenlijk 10 van 16 straffen verzameld.} \\
\texttt{Probeer het opnieuw!} \\
            }
        }
    }
    & & \\ \\

    \theutterance \stepcounter{utterance}  
    & & \multicolumn{4}{p{0.6\linewidth}}{
        \cellcolor[rgb]{0.9,0.9,0.9}{
            \makecell[{{p{\linewidth}}}]{
                \texttt{\tiny{[P1$\langle$GM]}}
                \texttt{Straf: (2,5) is niet leeg, maar bevat '│'.} \\
\\ 
\texttt{Jullie hebben gezamenlijk 10 van 16 straffen verzameld.} \\
\texttt{Probeer het opnieuw!} \\
            }
        }
    }
    & \\ \\

    \theutterance \stepcounter{utterance}  
    & \multicolumn{4}{p{0.6\linewidth}}{
        \cellcolor[rgb]{0.9,0.9,0.9}{
            \makecell[{{p{\linewidth}}}]{
                \texttt{\tiny{[P1$\rangle$GM]}}
                \texttt{BEWEEG: S, (7,4)} \\
            }
        }
    }
    & & \\ \\

    \theutterance \stepcounter{utterance}  
    & & & \multicolumn{2}{p{0.3\linewidth}}{
        \cellcolor[rgb]{0.9,0.9,0.9}{
            \makecell[{{p{\linewidth}}}]{
                \texttt{\tiny{[GM$|$GM]}}
                \texttt{BEWEEG: S, (7,4)} \\
            }
        }
    }
    & & \\ \\

    \theutterance \stepcounter{utterance}  
    & & & \multicolumn{2}{p{0.3\linewidth}}{
        \cellcolor[rgb]{0.9,0.9,0.9}{
            \makecell[{{p{\linewidth}}}]{
                \texttt{\tiny{[GM$|$GM]}}
                \texttt{Straf: (7,4) is niet leeg, maar bevat '─'.} \\
\\ 
\texttt{Jullie hebben gezamenlijk 11 van 16 straffen verzameld.} \\
\texttt{Probeer het opnieuw!} \\
            }
        }
    }
    & & \\ \\

    \theutterance \stepcounter{utterance}  
    & & \multicolumn{4}{p{0.6\linewidth}}{
        \cellcolor[rgb]{0.9,0.9,0.9}{
            \makecell[{{p{\linewidth}}}]{
                \texttt{\tiny{[P1$\langle$GM]}}
                \texttt{Straf: (7,4) is niet leeg, maar bevat '─'.} \\
\\ 
\texttt{Jullie hebben gezamenlijk 11 van 16 straffen verzameld.} \\
\texttt{Probeer het opnieuw!} \\
            }
        }
    }
    & \\ \\

    \theutterance \stepcounter{utterance}  
    & \multicolumn{4}{p{0.6\linewidth}}{
        \cellcolor[rgb]{0.9,0.9,0.9}{
            \makecell[{{p{\linewidth}}}]{
                \texttt{\tiny{[P1$\rangle$GM]}}
                \texttt{ZEG: Ik maak steeds fouten met de muren. Mijn raster: rij 5 heeft '│' op (3,5) en '│' op (7,5)? Nee wacht, laat ik mijn raster goed beschrijven. Randen: rij 1 kolom 5{-}6 is een dubbele muur, rij 2 kolom 5{-}6 dito. Rij 4 heeft muren op kolom 2{-}3 en 5{-}6. Rij 6 heeft muren op kolom 2{-}3 en 5{-}6. Rij 7 heeft muren op kolom 4 en 6. Mijn echt lege cellen: (1,1)(2,1)(3,1), (1,2)(2,2)(4,2)(5,2)(6,2)(7,2), (1,3)(2,3)(3,3)(4,3)(5,3)(6,3)(7,3), (1,4)(5,4)? eh nee... Kun jij jouw complete raster met muren beschrijven? Dan kan ik beter vergelijken.} \\
            }
        }
    }
    & & \\ \\

    \theutterance \stepcounter{utterance}  
    & & & \multicolumn{2}{p{0.3\linewidth}}{
        \cellcolor[rgb]{0.9,0.9,0.9}{
            \makecell[{{p{\linewidth}}}]{
                \texttt{\tiny{[GM$|$GM]}}
                \texttt{ZEG: Ik maak steeds fouten met de muren. Mijn raster: rij 5 heeft '│' op (3,5) en '│' op (7,5)? Nee wacht, laat ik mijn raster goed beschrijven. Randen: rij 1 kolom 5{-}6 is een dubbele muur, rij 2 kolom 5{-}6 dito. Rij 4 heeft muren op kolom 2{-}3 en 5{-}6. Rij 6 heeft muren op kolom 2{-}3 en 5{-}6. Rij 7 heeft muren op kolom 4 en 6. Mijn echt lege cellen: (1,1)(2,1)(3,1), (1,2)(2,2)(4,2)(5,2)(6,2)(7,2), (1,3)(2,3)(3,3)(4,3)(5,3)(6,3)(7,3), (1,4)(5,4)? eh nee... Kun jij jouw complete raster met muren beschrijven? Dan kan ik beter vergelijken.} \\
            }
        }
    }
    & & \\ \\

    \theutterance \stepcounter{utterance}  
    & & \multicolumn{4}{p{0.6\linewidth}}{
        \cellcolor[rgb]{0.9,0.9,0.9}{
            \makecell[{{p{\linewidth}}}]{
                \texttt{\tiny{[GM$\rangle$P2]}}
                \texttt{Je bericht is doorgestuurd naar de andere speler.} \\
\\ 
\\ 
\texttt{Je speelt momenteel ronde 17 van maximaal 28.} \\
\texttt{Jullie hebben gezamenlijk 11 van 16 straffen verzameld.} \\
\texttt{De andere speler stuurde dit bericht:} \\
\\ 
\\ 
\texttt{"Ik maak steeds fouten met de muren. Mijn raster: rij 5 heeft '│' op (3,5) en '│' op (7,5)? Nee wacht, laat ik mijn raster goed beschrijven. Randen: rij 1 kolom 5{-}6 is een dubbele muur, rij 2 kolom 5{-}6 dito. Rij 4 heeft muren op kolom 2{-}3 en 5{-}6. Rij 6 heeft muren op kolom 2{-}3 en 5{-}6. Rij 7 heeft muren op kolom 4 en 6. Mijn echt lege cellen: (1,1)(2,1)(3,1), (1,2)(2,2)(4,2)(5,2)(6,2)(7,2), (1,3)(2,3)(3,3)(4,3)(5,3)(6,3)(7,3), (1,4)(5,4)? eh nee... Kun jij jouw complete raster met muren beschrijven? Dan kan ik beter vergelijken."} \\
\\ 
\\ 
\texttt{Wat is je volgende opdracht?} \\
            }
        }
    }
    & \\ \\

    \theutterance \stepcounter{utterance}  
    & & & \multicolumn{4}{p{0.6\linewidth}}{
        \cellcolor[rgb]{0.9,0.9,0.9}{
            \makecell[{{p{\linewidth}}}]{
                \texttt{\tiny{[GM$\langle$P2]}}
                \texttt{ZEG: Mijn raster: Rij 1: muren op kolom 5{-}6 (╤ en ═). Rij 2: muur op kolom 5 (├). Rij 3: muur op kolom 5 (│). Rij 4: muren op kolom 2{-}3 (┌ en ─) en kolom 5{-}6 (┼ en ─). Rij 5: muren op kolom 2 (│) en kolom 5 (│). Rij 6: muren op kolom 1 (╟), 2{-}3 (─ en ┴), 4 (┬), 5{-}6 (─ en ┤). Rij 7: muren op kolom 4 (│) en kolom 6 (│). Mijn lege cellen: (1,1)(2,1)(3,1)(4,1)(7,1), (1,2)(2,2)(4,2)(6,2)(7,2), (1,3)(2,3)(3,3)(4,3)(5,3)(7,3), (2,4)(3,4)(4,4)(5,4)(6,4)(7,4), (1,5)(2,5)(4,5)(5,5)(6,5)(7,5), (1,6)(2,6)(3,6)(4,6)(5,6)(7,6), (1,7)(2,7)(4,7)(5,7)(6,7)(7,7). Gemeenschappelijk leeg met jouw lijst van eerder: nog steeds veel! Laten we bij S(2,5), A(5,5), D(4,5) blijven.} \\
            }
        }
    }
     \\ \\

    \theutterance \stepcounter{utterance}  
    & & & \multicolumn{2}{p{0.3\linewidth}}{
        \cellcolor[rgb]{0.9,0.9,0.9}{
            \makecell[{{p{\linewidth}}}]{
                \texttt{\tiny{[GM$|$GM]}}
                \texttt{ZEG: Mijn raster: Rij 1: muren op kolom 5{-}6 (╤ en ═). Rij 2: muur op kolom 5 (├). Rij 3: muur op kolom 5 (│). Rij 4: muren op kolom 2{-}3 (┌ en ─) en kolom 5{-}6 (┼ en ─). Rij 5: muren op kolom 2 (│) en kolom 5 (│). Rij 6: muren op kolom 1 (╟), 2{-}3 (─ en ┴), 4 (┬), 5{-}6 (─ en ┤). Rij 7: muren op kolom 4 (│) en kolom 6 (│). Mijn lege cellen: (1,1)(2,1)(3,1)(4,1)(7,1), (1,2)(2,2)(4,2)(6,2)(7,2), (1,3)(2,3)(3,3)(4,3)(5,3)(7,3), (2,4)(3,4)(4,4)(5,4)(6,4)(7,4), (1,5)(2,5)(4,5)(5,5)(6,5)(7,5), (1,6)(2,6)(3,6)(4,6)(5,6)(7,6), (1,7)(2,7)(4,7)(5,7)(6,7)(7,7). Gemeenschappelijk leeg met jouw lijst van eerder: nog steeds veel! Laten we bij S(2,5), A(5,5), D(4,5) blijven.} \\
            }
        }
    }
    & & \\ \\

    \theutterance \stepcounter{utterance}  
    & & \multicolumn{4}{p{0.6\linewidth}}{
        \cellcolor[rgb]{0.9,0.9,0.9}{
            \makecell[{{p{\linewidth}}}]{
                \texttt{\tiny{[P1$\langle$GM]}}
                \texttt{Je bericht is doorgestuurd naar de andere speler.} \\
\\ 
\\ 
\texttt{Je speelt momenteel ronde 17 van maximaal 28.} \\
\texttt{Jullie hebben gezamenlijk 11 van 16 straffen verzameld.} \\
\texttt{De andere speler stuurde dit bericht:} \\
\\ 
\\ 
\texttt{"Mijn raster: Rij 1: muren op kolom 5{-}6 (╤ en ═). Rij 2: muur op kolom 5 (├). Rij 3: muur op kolom 5 (│). Rij 4: muren op kolom 2{-}3 (┌ en ─) en kolom 5{-}6 (┼ en ─). Rij 5: muren op kolom 2 (│) en kolom 5 (│). Rij 6: muren op kolom 1 (╟), 2{-}3 (─ en ┴), 4 (┬), 5{-}6 (─ en ┤). Rij 7: muren op kolom 4 (│) en kolom 6 (│). Mijn lege cellen: (1,1)(2,1)(3,1)(4,1)(7,1), (1,2)(2,2)(4,2)(6,2)(7,2), (1,3)(2,3)(3,3)(4,3)(5,3)(7,3), (2,4)(3,4)(4,4)(5,4)(6,4)(7,4), (1,5)(2,5)(4,5)(5,5)(6,5)(7,5), (1,6)(2,6)(3,6)(4,6)(5,6)(7,6), (1,7)(2,7)(4,7)(5,7)(6,7)(7,7). Gemeenschappelijk leeg met jouw lijst van eerder: nog steeds veel! Laten we bij S(2,5), A(5,5), D(4,5) blijven."} \\
\\ 
\\ 
\texttt{Wat is je volgende opdracht?} \\
            }
        }
    }
    & \\ \\

    \theutterance \stepcounter{utterance}  
    & \multicolumn{4}{p{0.6\linewidth}}{
        \cellcolor[rgb]{0.9,0.9,0.9}{
            \makecell[{{p{\linewidth}}}]{
                \texttt{\tiny{[P1$\rangle$GM]}}
                \texttt{ZEG: Ah, ik zie het! Bij mij is (2,5) een muur '│', niet leeg. Mijn muur op (2,5) staat in jouw raster niet. Laten we S, A, D op posities zetten die echt gemeenschappelijk leeg zijn. Uit jouw lijst: gemeenschappelijk leeg zijn onder andere (1,5), (5,5), (6,5), (7,5) in rij 5 en (1,6), (2,6), (3,6), (4,6), (5,6), (7,6) in rij 6. Voorstel: S(5,5), A(6,5), D(4,5) {-} wacht, (4,5) is bij jou leeg. Is (4,5) bij mij leeg? Ja! Dus S(5,5), A(6,5), D(4,5). Mijn S staat op (1,4), A op (3,5), D op (3,7). Ik ga nu echt verplaatsen. Eerst S naar (5,5).} \\
            }
        }
    }
    & & \\ \\

    \theutterance \stepcounter{utterance}  
    & & & \multicolumn{2}{p{0.3\linewidth}}{
        \cellcolor[rgb]{0.9,0.9,0.9}{
            \makecell[{{p{\linewidth}}}]{
                \texttt{\tiny{[GM$|$GM]}}
                \texttt{ZEG: Ah, ik zie het! Bij mij is (2,5) een muur '│', niet leeg. Mijn muur op (2,5) staat in jouw raster niet. Laten we S, A, D op posities zetten die echt gemeenschappelijk leeg zijn. Uit jouw lijst: gemeenschappelijk leeg zijn onder andere (1,5), (5,5), (6,5), (7,5) in rij 5 en (1,6), (2,6), (3,6), (4,6), (5,6), (7,6) in rij 6. Voorstel: S(5,5), A(6,5), D(4,5) {-} wacht, (4,5) is bij jou leeg. Is (4,5) bij mij leeg? Ja! Dus S(5,5), A(6,5), D(4,5). Mijn S staat op (1,4), A op (3,5), D op (3,7). Ik ga nu echt verplaatsen. Eerst S naar (5,5).} \\
            }
        }
    }
    & & \\ \\

    \theutterance \stepcounter{utterance}  
    & & \multicolumn{4}{p{0.6\linewidth}}{
        \cellcolor[rgb]{0.9,0.9,0.9}{
            \makecell[{{p{\linewidth}}}]{
                \texttt{\tiny{[GM$\rangle$P2]}}
                \texttt{Je bericht is doorgestuurd naar de andere speler.} \\
\\ 
\\ 
\texttt{Je speelt momenteel ronde 18 van maximaal 28.} \\
\texttt{Jullie hebben gezamenlijk 11 van 16 straffen verzameld.} \\
\texttt{De andere speler stuurde dit bericht:} \\
\\ 
\\ 
\texttt{"Ah, ik zie het! Bij mij is (2,5) een muur '│', niet leeg. Mijn muur op (2,5) staat in jouw raster niet. Laten we S, A, D op posities zetten die echt gemeenschappelijk leeg zijn. Uit jouw lijst: gemeenschappelijk leeg zijn onder andere (1,5), (5,5), (6,5), (7,5) in rij 5 en (1,6), (2,6), (3,6), (4,6), (5,6), (7,6) in rij 6. Voorstel: S(5,5), A(6,5), D(4,5) {-} wacht, (4,5) is bij jou leeg. Is (4,5) bij mij leeg? Ja! Dus S(5,5), A(6,5), D(4,5). Mijn S staat op (1,4), A op (3,5), D op (3,7). Ik ga nu echt verplaatsen. Eerst S naar (5,5)."} \\
\\ 
\\ 
\texttt{Wat is je volgende opdracht?} \\
            }
        }
    }
    & \\ \\

    \theutterance \stepcounter{utterance}  
    & & & \multicolumn{4}{p{0.6\linewidth}}{
        \cellcolor[rgb]{0.9,0.9,0.9}{
            \makecell[{{p{\linewidth}}}]{
                \texttt{\tiny{[GM$\langle$P2]}}
                \texttt{ZEG: Wacht even, laten we eerst zeker weten dat (5,5) en (6,5) bij jou leeg zijn. Je zei eerder dat je muur op (2,5) hebt, klopt dat? En (4,5) is bij ons beide leeg. Laten we voor de zekerheid S(5,5), A(6,5), D(4,5) afspreken. Als jij bevestigt dat (5,5) en (6,5) leeg zijn bij jou, dan ga ik S naar (5,5) verplaatsen.} \\
            }
        }
    }
     \\ \\

    \theutterance \stepcounter{utterance}  
    & & & \multicolumn{2}{p{0.3\linewidth}}{
        \cellcolor[rgb]{0.9,0.9,0.9}{
            \makecell[{{p{\linewidth}}}]{
                \texttt{\tiny{[GM$|$GM]}}
                \texttt{ZEG: Wacht even, laten we eerst zeker weten dat (5,5) en (6,5) bij jou leeg zijn. Je zei eerder dat je muur op (2,5) hebt, klopt dat? En (4,5) is bij ons beide leeg. Laten we voor de zekerheid S(5,5), A(6,5), D(4,5) afspreken. Als jij bevestigt dat (5,5) en (6,5) leeg zijn bij jou, dan ga ik S naar (5,5) verplaatsen.} \\
            }
        }
    }
    & & \\ \\

    \theutterance \stepcounter{utterance}  
    & & \multicolumn{4}{p{0.6\linewidth}}{
        \cellcolor[rgb]{0.9,0.9,0.9}{
            \makecell[{{p{\linewidth}}}]{
                \texttt{\tiny{[P1$\langle$GM]}}
                \texttt{Je bericht is doorgestuurd naar de andere speler.} \\
\\ 
\\ 
\texttt{Je speelt momenteel ronde 18 van maximaal 28.} \\
\texttt{Jullie hebben gezamenlijk 11 van 16 straffen verzameld.} \\
\texttt{De andere speler stuurde dit bericht:} \\
\\ 
\\ 
\texttt{"Wacht even, laten we eerst zeker weten dat (5,5) en (6,5) bij jou leeg zijn. Je zei eerder dat je muur op (2,5) hebt, klopt dat? En (4,5) is bij ons beide leeg. Laten we voor de zekerheid S(5,5), A(6,5), D(4,5) afspreken. Als jij bevestigt dat (5,5) en (6,5) leeg zijn bij jou, dan ga ik S naar (5,5) verplaatsen."} \\
\\ 
\\ 
\texttt{Wat is je volgende opdracht?} \\
            }
        }
    }
    & \\ \\

    \theutterance \stepcounter{utterance}  
    & \multicolumn{4}{p{0.6\linewidth}}{
        \cellcolor[rgb]{0.9,0.9,0.9}{
            \makecell[{{p{\linewidth}}}]{
                \texttt{\tiny{[P1$\rangle$GM]}}
                \texttt{ZEG: Ja, (5,5) en (6,5) en (4,5) zijn allemaal leeg bij mij. Bevestigd: S(5,5), A(6,5), D(4,5). De rest staat al: M(2,2), P(3,1), O(4,1), T(3,3). Ik ga nu S verplaatsen!} \\
            }
        }
    }
    & & \\ \\

    \theutterance \stepcounter{utterance}  
    & & & \multicolumn{2}{p{0.3\linewidth}}{
        \cellcolor[rgb]{0.9,0.9,0.9}{
            \makecell[{{p{\linewidth}}}]{
                \texttt{\tiny{[GM$|$GM]}}
                \texttt{ZEG: Ja, (5,5) en (6,5) en (4,5) zijn allemaal leeg bij mij. Bevestigd: S(5,5), A(6,5), D(4,5). De rest staat al: M(2,2), P(3,1), O(4,1), T(3,3). Ik ga nu S verplaatsen!} \\
            }
        }
    }
    & & \\ \\

    \theutterance \stepcounter{utterance}  
    & & \multicolumn{4}{p{0.6\linewidth}}{
        \cellcolor[rgb]{0.9,0.9,0.9}{
            \makecell[{{p{\linewidth}}}]{
                \texttt{\tiny{[GM$\rangle$P2]}}
                \texttt{Je bericht is doorgestuurd naar de andere speler.} \\
\\ 
\\ 
\texttt{Je speelt momenteel ronde 19 van maximaal 28.} \\
\texttt{Jullie hebben gezamenlijk 11 van 16 straffen verzameld.} \\
\texttt{De andere speler stuurde dit bericht:} \\
\\ 
\\ 
\texttt{"Ja, (5,5) en (6,5) en (4,5) zijn allemaal leeg bij mij. Bevestigd: S(5,5), A(6,5), D(4,5). De rest staat al: M(2,2), P(3,1), O(4,1), T(3,3). Ik ga nu S verplaatsen!"} \\
\\ 
\\ 
\texttt{Wat is je volgende opdracht?} \\
            }
        }
    }
    & \\ \\

    \theutterance \stepcounter{utterance}  
    & & & \multicolumn{4}{p{0.6\linewidth}}{
        \cellcolor[rgb]{0.9,0.9,0.9}{
            \makecell[{{p{\linewidth}}}]{
                \texttt{\tiny{[GM$\langle$P2]}}
                \texttt{BEWEEG: S, (5,5)} \\
            }
        }
    }
     \\ \\

    \theutterance \stepcounter{utterance}  
    & & & \multicolumn{2}{p{0.3\linewidth}}{
        \cellcolor[rgb]{0.9,0.9,0.9}{
            \makecell[{{p{\linewidth}}}]{
                \texttt{\tiny{[GM$|$GM]}}
                \texttt{BEWEEG: S, (5,5)} \\
            }
        }
    }
    & & \\ \\

    \theutterance \stepcounter{utterance}  
    & & & \multicolumn{2}{p{0.3\linewidth}}{
        \cellcolor[rgb]{0.9,0.9,0.9}{
            \makecell[{{p{\linewidth}}}]{
                \texttt{\tiny{[GM$|$GM]}}
                \texttt{Straf: (5,5) is niet leeg, maar bevat '│'.} \\
\\ 
\texttt{Jullie hebben gezamenlijk 12 van 16 straffen verzameld.} \\
\texttt{Probeer het opnieuw!} \\
            }
        }
    }
    & & \\ \\

    \theutterance \stepcounter{utterance}  
    & & \multicolumn{4}{p{0.6\linewidth}}{
        \cellcolor[rgb]{0.9,0.9,0.9}{
            \makecell[{{p{\linewidth}}}]{
                \texttt{\tiny{[GM$\rangle$P2]}}
                \texttt{Straf: (5,5) is niet leeg, maar bevat '│'.} \\
\\ 
\texttt{Jullie hebben gezamenlijk 12 van 16 straffen verzameld.} \\
\texttt{Probeer het opnieuw!} \\
            }
        }
    }
    & \\ \\

    \theutterance \stepcounter{utterance}  
    & & & \multicolumn{4}{p{0.6\linewidth}}{
        \cellcolor[rgb]{0.9,0.9,0.9}{
            \makecell[{{p{\linewidth}}}]{
                \texttt{\tiny{[GM$\langle$P2]}}
                \texttt{BEWEEG: S, (6,5)} \\
            }
        }
    }
     \\ \\

    \theutterance \stepcounter{utterance}  
    & & & \multicolumn{2}{p{0.3\linewidth}}{
        \cellcolor[rgb]{0.9,0.9,0.9}{
            \makecell[{{p{\linewidth}}}]{
                \texttt{\tiny{[GM$|$GM]}}
                \texttt{BEWEEG: S, (6,5)} \\
            }
        }
    }
    & & \\ \\

    \theutterance \stepcounter{utterance}  
    & & & \multicolumn{2}{p{0.3\linewidth}}{
        \cellcolor[rgb]{0.9,0.9,0.9}{
            \makecell[{{p{\linewidth}}}]{
                \texttt{\tiny{[GM$|$GM]}}
                \texttt{S moved to (6, 5)} \\
            }
        }
    }
    & & \\ \\

    \theutterance \stepcounter{utterance}  
    & & \multicolumn{4}{p{0.6\linewidth}}{
        \cellcolor[rgb]{0.9,0.9,0.9}{
            \makecell[{{p{\linewidth}}}]{
                \texttt{\tiny{[P1$\langle$GM]}}
                \texttt{Je bericht is doorgestuurd naar de andere speler.} \\
\\ 
\\ 
\texttt{Je speelt momenteel ronde 19 van maximaal 28.} \\
\texttt{Jullie hebben gezamenlijk 12 van 16 straffen verzameld.} \\
\texttt{De andere speler heeft een voorwerp op zijn rooster verplaatst.} \\
\\ 
\\ 
\texttt{Wat is je volgende opdracht?} \\
            }
        }
    }
    & \\ \\

    \theutterance \stepcounter{utterance}  
    & \multicolumn{4}{p{0.6\linewidth}}{
        \cellcolor[rgb]{0.9,0.9,0.9}{
            \makecell[{{p{\linewidth}}}]{
                \texttt{\tiny{[P1$\rangle$GM]}}
                \texttt{BEWEEG: S, (5,5)} \\
            }
        }
    }
    & & \\ \\

    \theutterance \stepcounter{utterance}  
    & & & \multicolumn{2}{p{0.3\linewidth}}{
        \cellcolor[rgb]{0.9,0.9,0.9}{
            \makecell[{{p{\linewidth}}}]{
                \texttt{\tiny{[GM$|$GM]}}
                \texttt{BEWEEG: S, (5,5)} \\
            }
        }
    }
    & & \\ \\

    \theutterance \stepcounter{utterance}  
    & & & \multicolumn{2}{p{0.3\linewidth}}{
        \cellcolor[rgb]{0.9,0.9,0.9}{
            \makecell[{{p{\linewidth}}}]{
                \texttt{\tiny{[GM$|$GM]}}
                \texttt{Straf: (5,5) is niet leeg, maar bevat '│'.} \\
\\ 
\texttt{Jullie hebben gezamenlijk 13 van 16 straffen verzameld.} \\
\texttt{Probeer het opnieuw!} \\
            }
        }
    }
    & & \\ \\

    \theutterance \stepcounter{utterance}  
    & & \multicolumn{4}{p{0.6\linewidth}}{
        \cellcolor[rgb]{0.9,0.9,0.9}{
            \makecell[{{p{\linewidth}}}]{
                \texttt{\tiny{[P1$\langle$GM]}}
                \texttt{Straf: (5,5) is niet leeg, maar bevat '│'.} \\
\\ 
\texttt{Jullie hebben gezamenlijk 13 van 16 straffen verzameld.} \\
\texttt{Probeer het opnieuw!} \\
            }
        }
    }
    & \\ \\

    \theutterance \stepcounter{utterance}  
    & \multicolumn{4}{p{0.6\linewidth}}{
        \cellcolor[rgb]{0.9,0.9,0.9}{
            \makecell[{{p{\linewidth}}}]{
                \texttt{\tiny{[P1$\rangle$GM]}}
                \texttt{BEWEEG: S, (6,5)} \\
            }
        }
    }
    & & \\ \\

    \theutterance \stepcounter{utterance}  
    & & & \multicolumn{2}{p{0.3\linewidth}}{
        \cellcolor[rgb]{0.9,0.9,0.9}{
            \makecell[{{p{\linewidth}}}]{
                \texttt{\tiny{[GM$|$GM]}}
                \texttt{BEWEEG: S, (6,5)} \\
            }
        }
    }
    & & \\ \\

    \theutterance \stepcounter{utterance}  
    & & & \multicolumn{2}{p{0.3\linewidth}}{
        \cellcolor[rgb]{0.9,0.9,0.9}{
            \makecell[{{p{\linewidth}}}]{
                \texttt{\tiny{[GM$|$GM]}}
                \texttt{S moved to (6, 5)} \\
            }
        }
    }
    & & \\ \\

    \theutterance \stepcounter{utterance}  
    & & \multicolumn{4}{p{0.6\linewidth}}{
        \cellcolor[rgb]{0.9,0.9,0.9}{
            \makecell[{{p{\linewidth}}}]{
                \texttt{\tiny{[GM$\rangle$P2]}}
                \texttt{Verplaatst 'S' naar (6,5) met succes. Uw bijgewerkte raster ziet er als volgt uit:} \\
\texttt{```} \\
\texttt{ 1234567} \\
\texttt{╔════╤══╗} \\
\texttt{║◌◌PO│◌◌║ 1} \\
\texttt{║◌M◌◌├──╢ 2} \\
\texttt{║◌◌T◌│◌◌║ 3} \\
\texttt{║◌┌──┼──╢ 4} \\
\texttt{║A│◌◌│S◌║ 5} \\
\texttt{╟─┴┬─┤◌D║ 6} \\
\texttt{║◌◌│◌│◌◌║ 7} \\
\texttt{╚══╧═╧══╝} \\
\texttt{```} \\
\texttt{Je speelt momenteel ronde 20 van maximaal 28.} \\
\texttt{Jullie hebben gezamenlijk 13 van 16 straffen verzameld.} \\
\texttt{De andere speler heeft een voorwerp op zijn rooster verplaatst.} \\
\\ 
\\ 
\texttt{Wat is je volgende opdracht?} \\
            }
        }
    }
    & \\ \\

    \theutterance \stepcounter{utterance}  
    & & & \multicolumn{4}{p{0.6\linewidth}}{
        \cellcolor[rgb]{0.9,0.9,0.9}{
            \makecell[{{p{\linewidth}}}]{
                \texttt{\tiny{[GM$\langle$P2]}}
                \texttt{ZEG: Let op, (5,5) is bij mij een muur '│', niet leeg! Ik heb S naar (6,5) verplaatst. Mijn excuses, ik had dat niet goed gecheckt. Laten we de configuratie aanpassen. Mijn S staat nu op (6,5). Is (6,5) bij jou leeg? Anders moeten we S, A en D opnieuw afstemmen. Wat zijn nu jouw lege posities? Mijn overige: A moet nog naar een doel, D staat op (6,6).} \\
            }
        }
    }
     \\ \\

    \theutterance \stepcounter{utterance}  
    & & & \multicolumn{2}{p{0.3\linewidth}}{
        \cellcolor[rgb]{0.9,0.9,0.9}{
            \makecell[{{p{\linewidth}}}]{
                \texttt{\tiny{[GM$|$GM]}}
                \texttt{ZEG: Let op, (5,5) is bij mij een muur '│', niet leeg! Ik heb S naar (6,5) verplaatst. Mijn excuses, ik had dat niet goed gecheckt. Laten we de configuratie aanpassen. Mijn S staat nu op (6,5). Is (6,5) bij jou leeg? Anders moeten we S, A en D opnieuw afstemmen. Wat zijn nu jouw lege posities? Mijn overige: A moet nog naar een doel, D staat op (6,6).} \\
            }
        }
    }
    & & \\ \\

    \theutterance \stepcounter{utterance}  
    & & \multicolumn{4}{p{0.6\linewidth}}{
        \cellcolor[rgb]{0.9,0.9,0.9}{
            \makecell[{{p{\linewidth}}}]{
                \texttt{\tiny{[P1$\langle$GM]}}
                \texttt{Verplaatst 'S' naar (6,5) met succes. Uw bijgewerkte raster ziet er als volgt uit:} \\
\texttt{```} \\
\texttt{ 1234567} \\
\texttt{╔════╤══╗} \\
\texttt{║◌◌PO│◌◌║ 1} \\
\texttt{║◌M◌◌├──╢ 2} \\
\texttt{║◌◌T◌│◌◌║ 3} \\
\texttt{║◌┌──┼──╢ 4} \\
\texttt{║◌│◌A│S◌║ 5} \\
\texttt{╟─┴┬─┤◌◌║ 6} \\
\texttt{║◌◌│◌│D◌║ 7} \\
\texttt{╚══╧═╧══╝} \\
\texttt{```} \\
\texttt{Je speelt momenteel ronde 20 van maximaal 28.} \\
\texttt{Jullie hebben gezamenlijk 13 van 16 straffen verzameld.} \\
\texttt{De andere speler stuurde dit bericht:} \\
\\ 
\\ 
\texttt{"Let op, (5,5) is bij mij een muur '│', niet leeg! Ik heb S naar (6,5) verplaatst. Mijn excuses, ik had dat niet goed gecheckt. Laten we de configuratie aanpassen. Mijn S staat nu op (6,5). Is (6,5) bij jou leeg? Anders moeten we S, A en D opnieuw afstemmen. Wat zijn nu jouw lege posities? Mijn overige: A moet nog naar een doel, D staat op (6,6)."} \\
\\ 
\\ 
\texttt{Wat is je volgende opdracht?} \\
            }
        }
    }
    & \\ \\

    \theutterance \stepcounter{utterance}  
    & \multicolumn{4}{p{0.6\linewidth}}{
        \cellcolor[rgb]{0.9,0.9,0.9}{
            \makecell[{{p{\linewidth}}}]{
                \texttt{\tiny{[P1$\rangle$GM]}}
                \texttt{ZEG: Geen probleem, ik heb S ook op (6,5) gezet! Mijn A staat op (3,5) en D op (3,7). Laten we afspreken: S(6,5), A(1,5), D(4,5). (1,5) en (4,5) zijn bij mij leeg. Zijn die bij jou ook leeg?} \\
            }
        }
    }
    & & \\ \\

    \theutterance \stepcounter{utterance}  
    & & & \multicolumn{2}{p{0.3\linewidth}}{
        \cellcolor[rgb]{0.9,0.9,0.9}{
            \makecell[{{p{\linewidth}}}]{
                \texttt{\tiny{[GM$|$GM]}}
                \texttt{ZEG: Geen probleem, ik heb S ook op (6,5) gezet! Mijn A staat op (3,5) en D op (3,7). Laten we afspreken: S(6,5), A(1,5), D(4,5). (1,5) en (4,5) zijn bij mij leeg. Zijn die bij jou ook leeg?} \\
            }
        }
    }
    & & \\ \\

    \theutterance \stepcounter{utterance}  
    & & \multicolumn{4}{p{0.6\linewidth}}{
        \cellcolor[rgb]{0.9,0.9,0.9}{
            \makecell[{{p{\linewidth}}}]{
                \texttt{\tiny{[GM$\rangle$P2]}}
                \texttt{Je bericht is doorgestuurd naar de andere speler.} \\
\\ 
\\ 
\texttt{Je speelt momenteel ronde 21 van maximaal 28.} \\
\texttt{Jullie hebben gezamenlijk 13 van 16 straffen verzameld.} \\
\texttt{De andere speler stuurde dit bericht:} \\
\\ 
\\ 
\texttt{"Geen probleem, ik heb S ook op (6,5) gezet! Mijn A staat op (3,5) en D op (3,7). Laten we afspreken: S(6,5), A(1,5), D(4,5). (1,5) en (4,5) zijn bij mij leeg. Zijn die bij jou ook leeg?"} \\
\\ 
\\ 
\texttt{Wat is je volgende opdracht?} \\
            }
        }
    }
    & \\ \\

    \theutterance \stepcounter{utterance}  
    & & & \multicolumn{4}{p{0.6\linewidth}}{
        \cellcolor[rgb]{0.9,0.9,0.9}{
            \makecell[{{p{\linewidth}}}]{
                \texttt{\tiny{[GM$\langle$P2]}}
                \texttt{ZEG: Ja! (1,5) en (4,5) zijn bij mij leeg. Afgesproken: S(6,5), A(1,5), D(4,5). Ik ga A verplaatsen naar (1,5).} \\
            }
        }
    }
     \\ \\

    \theutterance \stepcounter{utterance}  
    & & & \multicolumn{2}{p{0.3\linewidth}}{
        \cellcolor[rgb]{0.9,0.9,0.9}{
            \makecell[{{p{\linewidth}}}]{
                \texttt{\tiny{[GM$|$GM]}}
                \texttt{ZEG: Ja! (1,5) en (4,5) zijn bij mij leeg. Afgesproken: S(6,5), A(1,5), D(4,5). Ik ga A verplaatsen naar (1,5).} \\
            }
        }
    }
    & & \\ \\

    \theutterance \stepcounter{utterance}  
    & & \multicolumn{4}{p{0.6\linewidth}}{
        \cellcolor[rgb]{0.9,0.9,0.9}{
            \makecell[{{p{\linewidth}}}]{
                \texttt{\tiny{[P1$\langle$GM]}}
                \texttt{Je bericht is doorgestuurd naar de andere speler.} \\
\\ 
\\ 
\texttt{Je speelt momenteel ronde 21 van maximaal 28.} \\
\texttt{Jullie hebben gezamenlijk 13 van 16 straffen verzameld.} \\
\texttt{De andere speler stuurde dit bericht:} \\
\\ 
\\ 
\texttt{"Ja! (1,5) en (4,5) zijn bij mij leeg. Afgesproken: S(6,5), A(1,5), D(4,5). Ik ga A verplaatsen naar (1,5)."} \\
\\ 
\\ 
\texttt{Wat is je volgende opdracht?} \\
            }
        }
    }
    & \\ \\

    \theutterance \stepcounter{utterance}  
    & \multicolumn{4}{p{0.6\linewidth}}{
        \cellcolor[rgb]{0.9,0.9,0.9}{
            \makecell[{{p{\linewidth}}}]{
                \texttt{\tiny{[P1$\rangle$GM]}}
                \texttt{BEWEEG: A, (1,5)} \\
            }
        }
    }
    & & \\ \\

    \theutterance \stepcounter{utterance}  
    & & & \multicolumn{2}{p{0.3\linewidth}}{
        \cellcolor[rgb]{0.9,0.9,0.9}{
            \makecell[{{p{\linewidth}}}]{
                \texttt{\tiny{[GM$|$GM]}}
                \texttt{BEWEEG: A, (1,5)} \\
            }
        }
    }
    & & \\ \\

    \theutterance \stepcounter{utterance}  
    & & & \multicolumn{2}{p{0.3\linewidth}}{
        \cellcolor[rgb]{0.9,0.9,0.9}{
            \makecell[{{p{\linewidth}}}]{
                \texttt{\tiny{[GM$|$GM]}}
                \texttt{A moved to (1, 5)} \\
            }
        }
    }
    & & \\ \\

    \theutterance \stepcounter{utterance}  
    & & \multicolumn{4}{p{0.6\linewidth}}{
        \cellcolor[rgb]{0.9,0.9,0.9}{
            \makecell[{{p{\linewidth}}}]{
                \texttt{\tiny{[GM$\rangle$P2]}}
                \texttt{Je bericht is doorgestuurd naar de andere speler.} \\
\\ 
\\ 
\texttt{Je speelt momenteel ronde 22 van maximaal 28.} \\
\texttt{Jullie hebben gezamenlijk 13 van 16 straffen verzameld.} \\
\texttt{De andere speler heeft een voorwerp op zijn rooster verplaatst.} \\
\\ 
\\ 
\texttt{Wat is je volgende opdracht?} \\
            }
        }
    }
    & \\ \\

    \theutterance \stepcounter{utterance}  
    & & & \multicolumn{4}{p{0.6\linewidth}}{
        \cellcolor[rgb]{0.9,0.9,0.9}{
            \makecell[{{p{\linewidth}}}]{
                \texttt{\tiny{[GM$\langle$P2]}}
                \texttt{BEWEEG: A, (1,5)} \\
            }
        }
    }
     \\ \\

    \theutterance \stepcounter{utterance}  
    & & & \multicolumn{2}{p{0.3\linewidth}}{
        \cellcolor[rgb]{0.9,0.9,0.9}{
            \makecell[{{p{\linewidth}}}]{
                \texttt{\tiny{[GM$|$GM]}}
                \texttt{BEWEEG: A, (1,5)} \\
            }
        }
    }
    & & \\ \\

    \theutterance \stepcounter{utterance}  
    & & & \multicolumn{2}{p{0.3\linewidth}}{
        \cellcolor[rgb]{0.9,0.9,0.9}{
            \makecell[{{p{\linewidth}}}]{
                \texttt{\tiny{[GM$|$GM]}}
                \texttt{Straf: (1,5) is niet leeg, maar bevat 'A'.} \\
\\ 
\texttt{Jullie hebben gezamenlijk 14 van 16 straffen verzameld.} \\
\texttt{Probeer het opnieuw!} \\
            }
        }
    }
    & & \\ \\

    \theutterance \stepcounter{utterance}  
    & & \multicolumn{4}{p{0.6\linewidth}}{
        \cellcolor[rgb]{0.9,0.9,0.9}{
            \makecell[{{p{\linewidth}}}]{
                \texttt{\tiny{[GM$\rangle$P2]}}
                \texttt{Straf: (1,5) is niet leeg, maar bevat 'A'.} \\
\\ 
\texttt{Jullie hebben gezamenlijk 14 van 16 straffen verzameld.} \\
\texttt{Probeer het opnieuw!} \\
            }
        }
    }
    & \\ \\

    \theutterance \stepcounter{utterance}  
    & & & \multicolumn{4}{p{0.6\linewidth}}{
        \cellcolor[rgb]{0.9,0.9,0.9}{
            \makecell[{{p{\linewidth}}}]{
                \texttt{\tiny{[GM$\langle$P2]}}
                \texttt{BEWEEG: A, (2,5)} \\
            }
        }
    }
     \\ \\

    \theutterance \stepcounter{utterance}  
    & & & \multicolumn{2}{p{0.3\linewidth}}{
        \cellcolor[rgb]{0.9,0.9,0.9}{
            \makecell[{{p{\linewidth}}}]{
                \texttt{\tiny{[GM$|$GM]}}
                \texttt{BEWEEG: A, (2,5)} \\
            }
        }
    }
    & & \\ \\

    \theutterance \stepcounter{utterance}  
    & & & \multicolumn{2}{p{0.3\linewidth}}{
        \cellcolor[rgb]{0.9,0.9,0.9}{
            \makecell[{{p{\linewidth}}}]{
                \texttt{\tiny{[GM$|$GM]}}
                \texttt{Straf: (2,5) is niet leeg, maar bevat '│'.} \\
\\ 
\texttt{Jullie hebben gezamenlijk 15 van 16 straffen verzameld.} \\
\texttt{Probeer het opnieuw!} \\
            }
        }
    }
    & & \\ \\

    \theutterance \stepcounter{utterance}  
    & & \multicolumn{4}{p{0.6\linewidth}}{
        \cellcolor[rgb]{0.9,0.9,0.9}{
            \makecell[{{p{\linewidth}}}]{
                \texttt{\tiny{[GM$\rangle$P2]}}
                \texttt{Straf: (2,5) is niet leeg, maar bevat '│'.} \\
\\ 
\texttt{Jullie hebben gezamenlijk 15 van 16 straffen verzameld.} \\
\texttt{Probeer het opnieuw!} \\
            }
        }
    }
    & \\ \\

    \theutterance \stepcounter{utterance}  
    & & & \multicolumn{4}{p{0.6\linewidth}}{
        \cellcolor[rgb]{0.9,0.9,0.9}{
            \makecell[{{p{\linewidth}}}]{
                \texttt{\tiny{[GM$\langle$P2]}}
                \texttt{ZEG: Help! Ik maak fouten. Mijn A staat al op (1,5) blijkbaar, en (2,5) is een muur '│'. Mijn raster na S{-}zet: A op (1,5), S op (6,5), D op (6,6). Mijn lege posities in rij 5 zijn nu (4,5) en (6,5) is bezet door S. Laten we A op (1,5) laten staan. Is (1,5) bij jou leeg? En D op (4,5), is die bij jou leeg?} \\
            }
        }
    }
     \\ \\

    \theutterance \stepcounter{utterance}  
    & & & \multicolumn{2}{p{0.3\linewidth}}{
        \cellcolor[rgb]{0.9,0.9,0.9}{
            \makecell[{{p{\linewidth}}}]{
                \texttt{\tiny{[GM$|$GM]}}
                \texttt{ZEG: Help! Ik maak fouten. Mijn A staat al op (1,5) blijkbaar, en (2,5) is een muur '│'. Mijn raster na S{-}zet: A op (1,5), S op (6,5), D op (6,6). Mijn lege posities in rij 5 zijn nu (4,5) en (6,5) is bezet door S. Laten we A op (1,5) laten staan. Is (1,5) bij jou leeg? En D op (4,5), is die bij jou leeg?} \\
            }
        }
    }
    & & \\ \\

    \theutterance \stepcounter{utterance}  
    & & \multicolumn{4}{p{0.6\linewidth}}{
        \cellcolor[rgb]{0.9,0.9,0.9}{
            \makecell[{{p{\linewidth}}}]{
                \texttt{\tiny{[P1$\langle$GM]}}
                \texttt{Verplaatst 'A' naar (1,5) met succes. Uw bijgewerkte raster ziet er als volgt uit:} \\
\texttt{```} \\
\texttt{ 1234567} \\
\texttt{╔════╤══╗} \\
\texttt{║◌◌PO│◌◌║ 1} \\
\texttt{║◌M◌◌├──╢ 2} \\
\texttt{║◌◌T◌│◌◌║ 3} \\
\texttt{║◌┌──┼──╢ 4} \\
\texttt{║A│◌◌│S◌║ 5} \\
\texttt{╟─┴┬─┤◌◌║ 6} \\
\texttt{║◌◌│◌│D◌║ 7} \\
\texttt{╚══╧═╧══╝} \\
\texttt{```} \\
\texttt{Je speelt momenteel ronde 22 van maximaal 28.} \\
\texttt{Jullie hebben gezamenlijk 15 van 16 straffen verzameld.} \\
\texttt{De andere speler stuurde dit bericht:} \\
\\ 
\\ 
\texttt{"Help! Ik maak fouten. Mijn A staat al op (1,5) blijkbaar, en (2,5) is een muur '│'. Mijn raster na S{-}zet: A op (1,5), S op (6,5), D op (6,6). Mijn lege posities in rij 5 zijn nu (4,5) en (6,5) is bezet door S. Laten we A op (1,5) laten staan. Is (1,5) bij jou leeg? En D op (4,5), is die bij jou leeg?"} \\
\\ 
\\ 
\texttt{Wat is je volgende opdracht?} \\
            }
        }
    }
    & \\ \\

    \theutterance \stepcounter{utterance}  
    & \multicolumn{4}{p{0.6\linewidth}}{
        \cellcolor[rgb]{0.9,0.9,0.9}{
            \makecell[{{p{\linewidth}}}]{
                \texttt{\tiny{[P1$\rangle$GM]}}
                \texttt{ZEG: Ja, A op (1,5) staat nu goed bij mij! En (4,5) is bij mij leeg. D op (4,5) is prima. Mijn D staat nog op (3,7). Ik ga D naar (4,5) verplaatsen. Dan zijn we bijna klaar! Alleen O moet nog naar (4,1), maar die staat al op (4,1) bij mij. Hoe staat het bij jou?} \\
            }
        }
    }
    & & \\ \\

    \theutterance \stepcounter{utterance}  
    & & & \multicolumn{2}{p{0.3\linewidth}}{
        \cellcolor[rgb]{0.9,0.9,0.9}{
            \makecell[{{p{\linewidth}}}]{
                \texttt{\tiny{[GM$|$GM]}}
                \texttt{ZEG: Ja, A op (1,5) staat nu goed bij mij! En (4,5) is bij mij leeg. D op (4,5) is prima. Mijn D staat nog op (3,7). Ik ga D naar (4,5) verplaatsen. Dan zijn we bijna klaar! Alleen O moet nog naar (4,1), maar die staat al op (4,1) bij mij. Hoe staat het bij jou?} \\
            }
        }
    }
    & & \\ \\

    \theutterance \stepcounter{utterance}  
    & & \multicolumn{4}{p{0.6\linewidth}}{
        \cellcolor[rgb]{0.9,0.9,0.9}{
            \makecell[{{p{\linewidth}}}]{
                \texttt{\tiny{[GM$\rangle$P2]}}
                \texttt{Je bericht is doorgestuurd naar de andere speler.} \\
\\ 
\\ 
\texttt{Je speelt momenteel ronde 23 van maximaal 28.} \\
\texttt{Jullie hebben gezamenlijk 15 van 16 straffen verzameld.} \\
\texttt{De andere speler stuurde dit bericht:} \\
\\ 
\\ 
\texttt{"Ja, A op (1,5) staat nu goed bij mij! En (4,5) is bij mij leeg. D op (4,5) is prima. Mijn D staat nog op (3,7). Ik ga D naar (4,5) verplaatsen. Dan zijn we bijna klaar! Alleen O moet nog naar (4,1), maar die staat al op (4,1) bij mij. Hoe staat het bij jou?"} \\
\\ 
\\ 
\texttt{Wat is je volgende opdracht?} \\
            }
        }
    }
    & \\ \\

    \theutterance \stepcounter{utterance}  
    & & & \multicolumn{4}{p{0.6\linewidth}}{
        \cellcolor[rgb]{0.9,0.9,0.9}{
            \makecell[{{p{\linewidth}}}]{
                \texttt{\tiny{[GM$\langle$P2]}}
                \texttt{BEWEEG: D, (4,5)} \\
            }
        }
    }
     \\ \\

    \theutterance \stepcounter{utterance}  
    & & & \multicolumn{2}{p{0.3\linewidth}}{
        \cellcolor[rgb]{0.9,0.9,0.9}{
            \makecell[{{p{\linewidth}}}]{
                \texttt{\tiny{[GM$|$GM]}}
                \texttt{BEWEEG: D, (4,5)} \\
            }
        }
    }
    & & \\ \\

    \theutterance \stepcounter{utterance}  
    & & & \multicolumn{2}{p{0.3\linewidth}}{
        \cellcolor[rgb]{0.9,0.9,0.9}{
            \makecell[{{p{\linewidth}}}]{
                \texttt{\tiny{[GM$|$GM]}}
                \texttt{D moved to (4, 5)} \\
            }
        }
    }
    & & \\ \\

    \theutterance \stepcounter{utterance}  
    & & \multicolumn{4}{p{0.6\linewidth}}{
        \cellcolor[rgb]{0.9,0.9,0.9}{
            \makecell[{{p{\linewidth}}}]{
                \texttt{\tiny{[P1$\langle$GM]}}
                \texttt{Je bericht is doorgestuurd naar de andere speler.} \\
\\ 
\\ 
\texttt{Je speelt momenteel ronde 23 van maximaal 28.} \\
\texttt{Jullie hebben gezamenlijk 15 van 16 straffen verzameld.} \\
\texttt{De andere speler heeft een voorwerp op zijn rooster verplaatst.} \\
\\ 
\\ 
\texttt{Wat is je volgende opdracht?} \\
            }
        }
    }
    & \\ \\

    \theutterance \stepcounter{utterance}  
    & \multicolumn{4}{p{0.6\linewidth}}{
        \cellcolor[rgb]{0.9,0.9,0.9}{
            \makecell[{{p{\linewidth}}}]{
                \texttt{\tiny{[P1$\rangle$GM]}}
                \texttt{BEWEEG: D, (4,5)} \\
            }
        }
    }
    & & \\ \\

    \theutterance \stepcounter{utterance}  
    & & & \multicolumn{2}{p{0.3\linewidth}}{
        \cellcolor[rgb]{0.9,0.9,0.9}{
            \makecell[{{p{\linewidth}}}]{
                \texttt{\tiny{[GM$|$GM]}}
                \texttt{BEWEEG: D, (4,5)} \\
            }
        }
    }
    & & \\ \\

    \theutterance \stepcounter{utterance}  
    & & & \multicolumn{2}{p{0.3\linewidth}}{
        \cellcolor[rgb]{0.9,0.9,0.9}{
            \makecell[{{p{\linewidth}}}]{
                \texttt{\tiny{[GM$|$GM]}}
                \texttt{D moved to (4, 5)} \\
            }
        }
    }
    & & \\ \\

    \theutterance \stepcounter{utterance}  
    & & \multicolumn{4}{p{0.6\linewidth}}{
        \cellcolor[rgb]{0.9,0.9,0.9}{
            \makecell[{{p{\linewidth}}}]{
                \texttt{\tiny{[GM$\rangle$P2]}}
                \texttt{Verplaatst 'D' naar (4,5) met succes. Uw bijgewerkte raster ziet er als volgt uit:} \\
\texttt{```} \\
\texttt{ 1234567} \\
\texttt{╔════╤══╗} \\
\texttt{║◌◌PO│◌◌║ 1} \\
\texttt{║◌M◌◌├──╢ 2} \\
\texttt{║◌◌T◌│◌◌║ 3} \\
\texttt{║◌┌──┼──╢ 4} \\
\texttt{║A│◌D│S◌║ 5} \\
\texttt{╟─┴┬─┤◌◌║ 6} \\
\texttt{║◌◌│◌│◌◌║ 7} \\
\texttt{╚══╧═╧══╝} \\
\texttt{```} \\
\texttt{Je speelt momenteel ronde 24 van maximaal 28.} \\
\texttt{Jullie hebben gezamenlijk 15 van 16 straffen verzameld.} \\
\texttt{De andere speler heeft een voorwerp op zijn rooster verplaatst.} \\
\\ 
\\ 
\texttt{Wat is je volgende opdracht?} \\
            }
        }
    }
    & \\ \\

    \theutterance \stepcounter{utterance}  
    & & & \multicolumn{4}{p{0.6\linewidth}}{
        \cellcolor[rgb]{0.9,0.9,0.9}{
            \makecell[{{p{\linewidth}}}]{
                \texttt{\tiny{[GM$\langle$P2]}}
                \texttt{ZEG: Mijn raster nu: M(2,2), P(3,1), O(4,1), T(3,3), A(1,5), D(4,5), S(6,5). Volgens afspraak: M(2,2), P(3,1), O(4,1), T(3,3), A(1,5), D(4,5), S(6,5). Alles staat goed! Is bij jou nu alles op de juiste plek? klaar?} \\
            }
        }
    }
     \\ \\

    \theutterance \stepcounter{utterance}  
    & & & \multicolumn{2}{p{0.3\linewidth}}{
        \cellcolor[rgb]{0.9,0.9,0.9}{
            \makecell[{{p{\linewidth}}}]{
                \texttt{\tiny{[GM$|$GM]}}
                \texttt{ZEG: Mijn raster nu: M(2,2), P(3,1), O(4,1), T(3,3), A(1,5), D(4,5), S(6,5). Volgens afspraak: M(2,2), P(3,1), O(4,1), T(3,3), A(1,5), D(4,5), S(6,5). Alles staat goed! Is bij jou nu alles op de juiste plek? klaar?} \\
            }
        }
    }
    & & \\ \\

    \theutterance \stepcounter{utterance}  
    & & \multicolumn{4}{p{0.6\linewidth}}{
        \cellcolor[rgb]{0.9,0.9,0.9}{
            \makecell[{{p{\linewidth}}}]{
                \texttt{\tiny{[P1$\langle$GM]}}
                \texttt{Verplaatst 'D' naar (4,5) met succes. Uw bijgewerkte raster ziet er als volgt uit:} \\
\texttt{```} \\
\texttt{ 1234567} \\
\texttt{╔════╤══╗} \\
\texttt{║◌◌PO│◌◌║ 1} \\
\texttt{║◌M◌◌├──╢ 2} \\
\texttt{║◌◌T◌│◌◌║ 3} \\
\texttt{║◌┌──┼──╢ 4} \\
\texttt{║A│◌D│S◌║ 5} \\
\texttt{╟─┴┬─┤◌◌║ 6} \\
\texttt{║◌◌│◌│◌◌║ 7} \\
\texttt{╚══╧═╧══╝} \\
\texttt{```} \\
\texttt{Je speelt momenteel ronde 24 van maximaal 28.} \\
\texttt{Jullie hebben gezamenlijk 15 van 16 straffen verzameld.} \\
\texttt{De andere speler stuurde dit bericht:} \\
\\ 
\\ 
\texttt{"Mijn raster nu: M(2,2), P(3,1), O(4,1), T(3,3), A(1,5), D(4,5), S(6,5). Volgens afspraak: M(2,2), P(3,1), O(4,1), T(3,3), A(1,5), D(4,5), S(6,5). Alles staat goed! Is bij jou nu alles op de juiste plek? klaar?"} \\
\\ 
\\ 
\texttt{Wat is je volgende opdracht?} \\
            }
        }
    }
    & \\ \\

    \theutterance \stepcounter{utterance}  
    & \multicolumn{4}{p{0.6\linewidth}}{
        \cellcolor[rgb]{0.9,0.9,0.9}{
            \makecell[{{p{\linewidth}}}]{
                \texttt{\tiny{[P1$\rangle$GM]}}
                \texttt{ZEG: Bij mij staat: M(2,2), P(3,1), O(4,1), T(3,3), S(6,5), A(1,5), D(4,5). Alles klopt! klaar!} \\
            }
        }
    }
    & & \\ \\

    \theutterance \stepcounter{utterance}  
    & & & \multicolumn{2}{p{0.3\linewidth}}{
        \cellcolor[rgb]{0.9,0.9,0.9}{
            \makecell[{{p{\linewidth}}}]{
                \texttt{\tiny{[GM$|$GM]}}
                \texttt{ZEG: Bij mij staat: M(2,2), P(3,1), O(4,1), T(3,3), S(6,5), A(1,5), D(4,5). Alles klopt! klaar!} \\
            }
        }
    }
    & & \\ \\

    \theutterance \stepcounter{utterance}  
    & & & \multicolumn{2}{p{0.3\linewidth}}{
        \cellcolor[rgb]{0.9,0.9,0.9}{
            \makecell[{{p{\linewidth}}}]{
                \texttt{\tiny{[GM$|$GM]}}
                \texttt{true} \\
            }
        }
    }
    & & \\ \\

    \theutterance \stepcounter{utterance}  
    & & & \multicolumn{2}{p{0.3\linewidth}}{
        \cellcolor[rgb]{0.9,0.9,0.9}{
            \makecell[{{p{\linewidth}}}]{
                \texttt{\tiny{[GM$|$GM]}}
                \texttt{Initial states:} \\
\texttt{```} \\
\texttt{Player 1:} \\
\texttt{ 1234567} \\
\texttt{╔════╤══╗} \\
\texttt{║◌◌◌O│PM║ 1} \\
\texttt{║◌◌T◌├──╢ 2} \\
\texttt{║◌◌◌◌│◌◌║ 3} \\
\texttt{║S┌──┼──╢ 4} \\
\texttt{║◌│◌A│◌◌║ 5} \\
\texttt{╟─┴┬─┤◌◌║ 6} \\
\texttt{║◌◌│◌│D◌║ 7} \\
\texttt{╚══╧═╧══╝} \\
\\ 
\\ 
\texttt{Player 2:} \\
\texttt{ 1234567} \\
\texttt{╔════╤══╗} \\
\texttt{║◌◌◌◌│O◌║ 1} \\
\texttt{║P◌M◌├──╢ 2} \\
\texttt{║◌S◌◌│T◌║ 3} \\
\texttt{║◌┌──┼──╢ 4} \\
\texttt{║A│◌◌│◌◌║ 5} \\
\texttt{╟─┴┬─┤◌D║ 6} \\
\texttt{║◌◌│◌│◌◌║ 7} \\
\texttt{╚══╧═╧══╝} \\
\\ 
\texttt{```} \\
            }
        }
    }
    & & \\ \\

    \theutterance \stepcounter{utterance}  
    & & & \multicolumn{2}{p{0.3\linewidth}}{
        \cellcolor[rgb]{0.9,0.9,0.9}{
            \makecell[{{p{\linewidth}}}]{
                \texttt{\tiny{[GM$|$GM]}}
                \texttt{End states:} \\
\texttt{```} \\
\texttt{Player 1:} \\
\texttt{ 1234567} \\
\texttt{╔════╤══╗} \\
\texttt{║◌◌PO│◌◌║ 1} \\
\texttt{║◌M◌◌├──╢ 2} \\
\texttt{║◌◌T◌│◌◌║ 3} \\
\texttt{║◌┌──┼──╢ 4} \\
\texttt{║A│◌D│S◌║ 5} \\
\texttt{╟─┴┬─┤◌◌║ 6} \\
\texttt{║◌◌│◌│◌◌║ 7} \\
\texttt{╚══╧═╧══╝} \\
\\ 
\\ 
\texttt{Player 2:} \\
\texttt{ 1234567} \\
\texttt{╔════╤══╗} \\
\texttt{║◌◌PO│◌◌║ 1} \\
\texttt{║◌M◌◌├──╢ 2} \\
\texttt{║◌◌T◌│◌◌║ 3} \\
\texttt{║◌┌──┼──╢ 4} \\
\texttt{║A│◌D│S◌║ 5} \\
\texttt{╟─┴┬─┤◌◌║ 6} \\
\texttt{║◌◌│◌│◌◌║ 7} \\
\texttt{╚══╧═╧══╝} \\
\\ 
\texttt{```} \\
            }
        }
    }
    & & \\ \\

    \theutterance \stepcounter{utterance}  
    & & & \multicolumn{2}{p{0.3\linewidth}}{
        \cellcolor[rgb]{0.9,0.9,0.9}{
            \makecell[{{p{\linewidth}}}]{
                \texttt{\tiny{[GM$|$GM]}}
                \texttt{* success: True} \\
\texttt{* lose: False} \\
\texttt{* aborted: False} \\
\texttt{{-}{-}{-}{-}{-}{-}{-}} \\
\texttt{* Shifts: 6.00} \\
\texttt{* Max Shifts: 12.00} \\
\texttt{* Min Shifts: 6.00} \\
\texttt{* End Distance Sum: 0.00} \\
\texttt{* Init Distance Sum: 20.21} \\
\texttt{* Expected Distance Sum: 29.33} \\
\texttt{* Penalties: 15.00} \\
\texttt{* Max Penalties: 16.00} \\
\texttt{* Rounds: 24.00} \\
\texttt{* Max Rounds: 28.00} \\
\texttt{* Object Count: 7.00} \\
            }
        }
    }
    & & \\ \\

\end{supertabular}
}

\end{document}
